\documentclass[12pt,a4paper]{article}
\usepackage[utf8]{inputenc}
\usepackage[T1]{fontenc}
\usepackage[brazilian]{babel}
\usepackage{geometry}
\usepackage{booktabs}
\usepackage{graphicx}
\usepackage{amsmath}
\usepackage{float}
\usepackage{longtable}
\geometry{margin=2.5cm}
\newcommand{\TODO}[1]{\textbf{[TODO: #1]}}
\newcommand{\grafico}[2]{%
  \IfFileExists{../experimentos/graficos/#1.pdf}{%
    \begin{figure}[H]\centering
    \includegraphics[width=0.9\textwidth]{../experimentos/graficos/#1.pdf}
    \caption{#2}\end{figure}%
  }{\TODO{gerar o gráfico #1}}%
}
\title{Projeto 2: Índices em Árvores de Busca}
\author{SCC0202 -- Algoritmos e Estruturas de Dados I}
\date{2026}

\begin{document}
\maketitle

\section{Objetivo e escolhas de implementação}
O programa indexa os CSVs do Projeto 1 sem armazenar a base inteira. O leitor devolve um registro de cada vez e reutiliza o buffer de campos. A árvore mantém uma chave textual por nó e um vetor de offsets em bytes para todas as linhas com aquele valor. O leitor ignora linhas totalmente vazias, preserva valores ausentes como strings vazias e entende aspas, delimitadores citados e quebras de linha citadas.

ABB, AVL e rubro-negra são TADs distintos. A ABB pode ficar degenerada, mas suas operações de percurso, leitura, gravação e destruição são iterativas para não estourar a pilha. A AVL grava altura por nó; a rubro-negra grava a cor de cada nó e calcula a altura ao final. As três mantêm um único nó por chave distinta.

As chaves são strings e as consultas usam igualdade exata com \texttt{strcmp}. Não há conversão numérica: \texttt{2020}, \texttt{2020.0} e \texttt{02020} são chaves diferentes. Colunas multivaloradas, como \texttt{Genres}, são indexadas como o texto completo; \texttt{Action} não equivale a \texttt{Action,Indie}.

\section{Formato persistido do índice}
O índice usa a assinatura de oito bytes \texttt{AED1IDX2}, seguida dos campos abaixo. Os inteiros de 32 e 64 bits são serializados byte a byte em little-endian. Nenhum ponteiro de memória é persistido.
\begin{center}
\begin{tabular}{ll}
\toprule
Campo & Representação \\
\midrule
Versão do formato & $u32$, atualmente 3 \\
Tipo de árvore & $u32$: ABB, AVL ou rubro-negra \\
Tipo de chave & $u32$: 1 para string \\
Nome da chave & comprimento $u32$ e bytes \\
Nome-base do CSV & comprimento $u32$ e bytes \\
Tamanho do CSV & $u64$ bytes \\
Modificação do CSV & segundos desde a época, armazenados em $u64$ \\
Registros indexados & $u64$ \\
Chaves distintas & $u64$ \\
Checksum do CSV & FNV-1a de 64 bits \\
\bottomrule
\end{tabular}
\end{center}

O checksum combina o cabeçalho CSV, os primeiros 4096 bytes e os últimos 4096 bytes. Na carga, o programa valida nome-base, tamanho, data de modificação e checksum; depois valida offsets, contagem de registros, número de chaves distintas, ordenação da árvore e os invariantes de balanceamento aplicáveis. Cada nó guarda um marcador de presença, a chave e a coleção de offsets. AVL também guarda sua altura em $u32$; rubro-negra guarda sua cor em um byte. Filhos são serializados em pré-ordem com marcadores de presença.

\section{Convenções e métricas}
A altura é medida em nós: árvore vazia tem altura zero e folha tem altura um. Uma rotação dupla soma duas rotações elementares. Recolors contam apenas mudanças efetivas; a cor inicial do nó recém-criado não é contada. \texttt{INSERT\_COMPARISONS} soma comparações durante inserções. \texttt{SEARCH\_COMPARISONS} é contado localmente durante buscas e não altera o TAD.

\texttt{MATCHED\_OFFSETS} soma os tamanhos dos vetores retornados por todas as linhas de consulta, incluindo chaves repetidas e sobreposição entre índices. \texttt{OUTPUT\_RECORDS} conta offsets únicos exportados. Os vetores de offsets de chaves consultadas são deduplicados por identidade antes de cópia; a união final é ordenada por offset e deduplicada. \texttt{SEARCH\_MS} cobre somente chamadas à árvore. \texttt{DEDUP\_MS}, \texttt{FETCH\_MS} e \texttt{WRITE\_MS} separam união, leitura direta do CSV e gravação. \texttt{READ\_MS} e \texttt{INSERT\_MS} são medidos por registro, portanto incluem o pequeno custo das chamadas ao relógio.

\section{Complexidade}
Se $D$ é o número de chaves distintas e $N$ o número de registros, a busca e inserção na ABB custam $O(h)$, em que $h$ é a altura. Com ordem de inserção aleatória, a altura esperada é $O(\log D)$; no pior caso, chaves em ordem crescente ou decrescente formam uma cadeia, com custo $O(D)$. Construir a ABB custa $O(N\log D)$ em ordem de inserção favorável, mas chega a $O(ND)$, ou $O(N^2)$ quando todas as chaves são distintas e inseridas em ordem.

AVL e rubro-negra mantêm altura $O(\log D)$ no pior caso. Busca e inserção custam $O(\log D)$ por registro consultado/inserido e a construção custa $O(N\log D)$. Uma chave repetida não cria nó nem rotação; a busca para localizar o nó continua custando $O(\log D)$ e anexar o offset custa $O(1)$ amortizado no vetor dinâmico. O armazenamento é $O(D+N)$, contando os nós e todos os offsets.

Para $Q$ chaves consultadas, a busca nas árvores custa $O(Qh)$ no total. Se as chaves encontradas apontam para $M$ offsets, mesclar e ordenar os resultados custa $O(M\log M)$ no pior caso. A aplicação busca registros por posicionamento no offset: cada acesso envolve um reposicionamento no arquivo e a leitura dos bytes do registro, seguida da escrita desses bytes. O custo físico de seeks e do volume transferido não é representado completamente pela complexidade assintótica.

\section{Discussão das estruturas}
\subsection{Quando a ABB se aproxima das balanceadas?}
Quando a ordem de inserção distribui as chaves de modo que os subárvores fiquem aproximadamente equilibrados, a ABB pode ter altura semelhante à AVL e à rubro-negra. Se as chaves chegam ordenadas, a árvore se aproxima de uma lista e tanto inserção quanto busca podem custar linearmente. A ABB iterativa evita falha de pilha nesse caso, sem remover a degradação de tempo esperada.

\subsection{Efeito de rotações e recolors}
Rotações e alterações de cor acrescentam operações durante a construção. A AVL atualiza alturas e pode executar rotações; a rubro-negra faz reparos de cor e rotações. Essas operações mantêm a altura limitada e, portanto, não são repetidas durante busca. O custo de busca depende da altura final e das comparações, não do número histórico de rotações ou recolors.

\subsection{Mais chaves simples ou complexas}
Em árvores balanceadas, aumentar o número de chaves consultadas aumenta aproximadamente de forma linear o número de buscas, com custo $O(Q\log D)$, além dos resultados retornados. Uma consulta complexa continua sendo uma união OU; ela usa os índices de cada coluna mencionada. O número total de chaves e a distribuição por coluna determinam o trabalho das árvores, enquanto o número de offsets retornados determina o custo posterior de mesclagem e acesso ao arquivo.

\subsection{Quando o arquivo domina?}
Quando muitas consultas retornam muitos registros, o custo de reposicionar o arquivo, transferir registros grandes e gravar a saída pode superar as comparações da árvore. Deduplicar por offset evita reler e regravar o mesmo registro quando consultas repetidas ou índices diferentes apontam para ele. Em discos com alto custo de seek, a ordem crescente dos offsets também pode reduzir movimentos em comparação a uma ordem arbitrária.

\subsection{Consistência entre repetições e bases}
Os tempos foram medidos em dez repetições e a ordem das árvores foi rotacionada. A árvore mais rápida variou entre base, coluna e tamanho; não há vencedora uniforme. Por exemplo, nas amostras de 10000 registros de asteroides, a rubro-negra teve a menor média de construção nas chaves \texttt{year} e \texttt{on\_sentry\_list}; nas amostras de 1000 registros, a AVL teve a menor média nas duas colunas de games. Em games, as diferenças médias de construção nas duas colunas e nas buscas de 1000 chaves também foram pequenas frente à variação de algumas configurações. A variabilidade é visível nos desvios padrão das tabelas, sobretudo na amostra de games com 1000 registros.

O aumento de chaves consultadas eleva o trabalho das árvores em ambas as modalidades. A consulta complexa busca em mais de um índice e une os offsets, por isso seu custo total também depende da sobreposição e da quantidade de registros encontrados. Nas configurações de 10000 registros e 1000 chaves, os tempos de árvore ficaram em décimos de milissegundo, enquanto o tempo total chegou a dezenas de milissegundos. Nesses casos, leitura dos registros e escrita da saída dominam. Remover offsets repetidos evita reler e regravar registros sobrepostos; a métrica \texttt{MATCHED\_OFFSETS} registra a quantidade bruta e \texttt{OUTPUT\_RECORDS}, a quantidade efetivamente exportada.

Essas observações são limitadas a dois tamanhos pequenos e a uma máquina. Frequência dinâmica do processador, cache, volume da saída e carga do sistema afetam as medições. A ausência de diferença estatística nos outros três pares testados não prova equivalência entre as estruturas; indica somente que, nessas configurações e com dez repetições, não se detectou diferença pelo teste usado.

\section{Protocolo experimental e resultados}
As medições deste relatório foram feitas nas duas bases locais do Projeto 1. O script seleciona registros por reservoir sampling, embaralha a amostra e fixa a semente em 20261125. Foram usados tamanhos de 1000 e 10000 registros, dez repetições por configuração, buscas simples com 1, 10, 100 e 1000 chaves e buscas complexas com 10, 100 e 1000 chaves distribuídas entre duas colunas. As consultas incluem valores presentes, inexistentes e repetidos. A ordem das árvores é rotacionada entre repetições. As saídas foram comparadas por MD5 em cada configuração; todas coincidiram.

O escopo medido é menor que os tamanhos máximos oferecidos pelos scripts (100000 e 500000 registros). Portanto, as observações abaixo descrevem apenas as amostras de 1000 e 10000 registros e não devem ser extrapoladas para bases maiores. Os CSVs brutos, o ambiente de execução e os resumos estão em \texttt{experimentos/dados/}; os gráficos são gerados em \texttt{experimentos/graficos/}.

O experimento foi executado em Linux x86\_64, kernel 6.18.62-1-lts, GCC 16.2.1 e flags \texttt{-std=c99 -Wall -Wextra -Wpedantic -O2}, em Intel Core i5-10210U. A coleta foi feita em uma única máquina; variação de frequência, cache e carga concorrente continuam sendo possíveis fontes de ruído.

\subsection{Construção}
\IfFileExists{tabela_construcao.tex}{\scriptsize
\begin{longtable}{p{4cm}p{1.4cm}p{2.6cm}rr}
\hline
Base/configuração & Árvore & Métrica & Média & Desvio\\
\hline
\endfirsthead
\hline
Base/configuração & Árvore & Métrica & Média & Desvio\\
\hline
\endhead
asteroids, 1000, on\_\allowbreak\{\}sentry\_\allowbreak\{\}list, ,  & abb & BUILD\_\allowbreak\{\}MS & 4.1944 & 2.15137\\
asteroids, 1000, on\_\allowbreak\{\}sentry\_\allowbreak\{\}list, ,  & abb & INSERT\_\allowbreak\{\}COMPARISONS & 1060 & 0\\
asteroids, 1000, on\_\allowbreak\{\}sentry\_\allowbreak\{\}list, ,  & abb & RECOLORS & 0 & 0\\
asteroids, 1000, on\_\allowbreak\{\}sentry\_\allowbreak\{\}list, ,  & abb & ROTATIONS & 0 & 0\\
asteroids, 1000, on\_\allowbreak\{\}sentry\_\allowbreak\{\}list, ,  & abb & TREE\_\allowbreak\{\}HEIGHT & 2 & 0\\
asteroids, 1000, year, ,  & abb & BUILD\_\allowbreak\{\}MS & 4.5529 & 2.28804\\
asteroids, 1000, year, ,  & abb & INSERT\_\allowbreak\{\}COMPARISONS & 5813 & 0\\
asteroids, 1000, year, ,  & abb & RECOLORS & 0 & 0\\
asteroids, 1000, year, ,  & abb & ROTATIONS & 0 & 0\\
asteroids, 1000, year, ,  & abb & TREE\_\allowbreak\{\}HEIGHT & 12 & 0\\
asteroids, 1000, on\_\allowbreak\{\}sentry\_\allowbreak\{\}list, ,  & avl & BUILD\_\allowbreak\{\}MS & 5.0066 & 2.62843\\
asteroids, 1000, on\_\allowbreak\{\}sentry\_\allowbreak\{\}list, ,  & avl & INSERT\_\allowbreak\{\}COMPARISONS & 1060 & 0\\
asteroids, 1000, on\_\allowbreak\{\}sentry\_\allowbreak\{\}list, ,  & avl & RECOLORS & 0 & 0\\
asteroids, 1000, on\_\allowbreak\{\}sentry\_\allowbreak\{\}list, ,  & avl & ROTATIONS & 0 & 0\\
asteroids, 1000, on\_\allowbreak\{\}sentry\_\allowbreak\{\}list, ,  & avl & TREE\_\allowbreak\{\}HEIGHT & 2 & 0\\
asteroids, 1000, year, ,  & avl & BUILD\_\allowbreak\{\}MS & 5.3203 & 2.87834\\
asteroids, 1000, year, ,  & avl & INSERT\_\allowbreak\{\}COMPARISONS & 5363 & 0\\
asteroids, 1000, year, ,  & avl & RECOLORS & 0 & 0\\
asteroids, 1000, year, ,  & avl & ROTATIONS & 53 & 0\\
asteroids, 1000, year, ,  & avl & TREE\_\allowbreak\{\}HEIGHT & 8 & 0\\
asteroids, 1000, on\_\allowbreak\{\}sentry\_\allowbreak\{\}list, ,  & rubro\_\allowbreak\{\}negra & BUILD\_\allowbreak\{\}MS & 4.3248 & 1.91306\\
asteroids, 1000, on\_\allowbreak\{\}sentry\_\allowbreak\{\}list, ,  & rubro\_\allowbreak\{\}negra & INSERT\_\allowbreak\{\}COMPARISONS & 1060 & 0\\
asteroids, 1000, on\_\allowbreak\{\}sentry\_\allowbreak\{\}list, ,  & rubro\_\allowbreak\{\}negra & RECOLORS & 1 & 0\\
asteroids, 1000, on\_\allowbreak\{\}sentry\_\allowbreak\{\}list, ,  & rubro\_\allowbreak\{\}negra & ROTATIONS & 0 & 0\\
asteroids, 1000, on\_\allowbreak\{\}sentry\_\allowbreak\{\}list, ,  & rubro\_\allowbreak\{\}negra & TREE\_\allowbreak\{\}HEIGHT & 2 & 0\\
asteroids, 1000, year, ,  & rubro\_\allowbreak\{\}negra & BUILD\_\allowbreak\{\}MS & 4.8369 & 2.0901\\
asteroids, 1000, year, ,  & rubro\_\allowbreak\{\}negra & INSERT\_\allowbreak\{\}COMPARISONS & 5446 & 0\\
asteroids, 1000, year, ,  & rubro\_\allowbreak\{\}negra & RECOLORS & 187 & 0\\
asteroids, 1000, year, ,  & rubro\_\allowbreak\{\}negra & ROTATIONS & 43 & 0\\
asteroids, 1000, year, ,  & rubro\_\allowbreak\{\}negra & TREE\_\allowbreak\{\}HEIGHT & 8 & 0\\
asteroids, 10000, on\_\allowbreak\{\}sentry\_\allowbreak\{\}list, ,  & abb & BUILD\_\allowbreak\{\}MS & 58.8758 & 13.1312\\
asteroids, 10000, on\_\allowbreak\{\}sentry\_\allowbreak\{\}list, ,  & abb & INSERT\_\allowbreak\{\}COMPARISONS & 10594 & 0\\
asteroids, 10000, on\_\allowbreak\{\}sentry\_\allowbreak\{\}list, ,  & abb & RECOLORS & 0 & 0\\
asteroids, 10000, on\_\allowbreak\{\}sentry\_\allowbreak\{\}list, ,  & abb & ROTATIONS & 0 & 0\\
asteroids, 10000, on\_\allowbreak\{\}sentry\_\allowbreak\{\}list, ,  & abb & TREE\_\allowbreak\{\}HEIGHT & 2 & 0\\
asteroids, 10000, year, ,  & abb & BUILD\_\allowbreak\{\}MS & 62.0717 & 16.8126\\
asteroids, 10000, year, ,  & abb & INSERT\_\allowbreak\{\}COMPARISONS & 69530 & 0\\
asteroids, 10000, year, ,  & abb & RECOLORS & 0 & 0\\
asteroids, 10000, year, ,  & abb & ROTATIONS & 0 & 0\\
asteroids, 10000, year, ,  & abb & TREE\_\allowbreak\{\}HEIGHT & 13 & 0\\
asteroids, 10000, on\_\allowbreak\{\}sentry\_\allowbreak\{\}list, ,  & avl & BUILD\_\allowbreak\{\}MS & 58.417 & 11.387\\
asteroids, 10000, on\_\allowbreak\{\}sentry\_\allowbreak\{\}list, ,  & avl & INSERT\_\allowbreak\{\}COMPARISONS & 10594 & 0\\
asteroids, 10000, on\_\allowbreak\{\}sentry\_\allowbreak\{\}list, ,  & avl & RECOLORS & 0 & 0\\
asteroids, 10000, on\_\allowbreak\{\}sentry\_\allowbreak\{\}list, ,  & avl & ROTATIONS & 0 & 0\\
asteroids, 10000, on\_\allowbreak\{\}sentry\_\allowbreak\{\}list, ,  & avl & TREE\_\allowbreak\{\}HEIGHT & 2 & 0\\
asteroids, 10000, year, ,  & avl & BUILD\_\allowbreak\{\}MS & 63.5089 & 9.08893\\
asteroids, 10000, year, ,  & avl & INSERT\_\allowbreak\{\}COMPARISONS & 55400 & 0\\
asteroids, 10000, year, ,  & avl & RECOLORS & 0 & 0\\
asteroids, 10000, year, ,  & avl & ROTATIONS & 53 & 0\\
asteroids, 10000, year, ,  & avl & TREE\_\allowbreak\{\}HEIGHT & 7 & 0\\
asteroids, 10000, on\_\allowbreak\{\}sentry\_\allowbreak\{\}list, ,  & rubro\_\allowbreak\{\}negra & BUILD\_\allowbreak\{\}MS & 57.3639 & 11.5431\\
asteroids, 10000, on\_\allowbreak\{\}sentry\_\allowbreak\{\}list, ,  & rubro\_\allowbreak\{\}negra & INSERT\_\allowbreak\{\}COMPARISONS & 10594 & 0\\
asteroids, 10000, on\_\allowbreak\{\}sentry\_\allowbreak\{\}list, ,  & rubro\_\allowbreak\{\}negra & RECOLORS & 1 & 0\\
asteroids, 10000, on\_\allowbreak\{\}sentry\_\allowbreak\{\}list, ,  & rubro\_\allowbreak\{\}negra & ROTATIONS & 0 & 0\\
asteroids, 10000, on\_\allowbreak\{\}sentry\_\allowbreak\{\}list, ,  & rubro\_\allowbreak\{\}negra & TREE\_\allowbreak\{\}HEIGHT & 2 & 0\\
asteroids, 10000, year, ,  & rubro\_\allowbreak\{\}negra & BUILD\_\allowbreak\{\}MS & 61.8322 & 15.8732\\
asteroids, 10000, year, ,  & rubro\_\allowbreak\{\}negra & INSERT\_\allowbreak\{\}COMPARISONS & 54948 & 0\\
asteroids, 10000, year, ,  & rubro\_\allowbreak\{\}negra & RECOLORS & 167 & 0\\
asteroids, 10000, year, ,  & rubro\_\allowbreak\{\}negra & ROTATIONS & 44 & 0\\
asteroids, 10000, year, ,  & rubro\_\allowbreak\{\}negra & TREE\_\allowbreak\{\}HEIGHT & 8 & 0\\
games, 1000, Genres, ,  & abb & BUILD\_\allowbreak\{\}MS & 56.6341 & 28.743\\
games, 1000, Genres, ,  & abb & INSERT\_\allowbreak\{\}COMPARISONS & 7667 & 0\\
games, 1000, Genres, ,  & abb & RECOLORS & 0 & 0\\
games, 1000, Genres, ,  & abb & ROTATIONS & 0 & 0\\
games, 1000, Genres, ,  & abb & TREE\_\allowbreak\{\}HEIGHT & 17 & 0\\
games, 1000, Release date, ,  & abb & BUILD\_\allowbreak\{\}MS & 54.8429 & 26.1079\\
games, 1000, Release date, ,  & abb & INSERT\_\allowbreak\{\}COMPARISONS & 13133 & 0\\
games, 1000, Release date, ,  & abb & RECOLORS & 0 & 0\\
games, 1000, Release date, ,  & abb & ROTATIONS & 0 & 0\\
games, 1000, Release date, ,  & abb & TREE\_\allowbreak\{\}HEIGHT & 26 & 0\\
games, 1000, Genres, ,  & avl & BUILD\_\allowbreak\{\}MS & 49.9994 & 24.4058\\
games, 1000, Genres, ,  & avl & INSERT\_\allowbreak\{\}COMPARISONS & 6293 & 0\\
games, 1000, Genres, ,  & avl & RECOLORS & 0 & 0\\
games, 1000, Genres, ,  & avl & ROTATIONS & 181 & 0\\
games, 1000, Genres, ,  & avl & TREE\_\allowbreak\{\}HEIGHT & 10 & 0\\
games, 1000, Release date, ,  & avl & BUILD\_\allowbreak\{\}MS & 54.5924 & 27.726\\
games, 1000, Release date, ,  & avl & INSERT\_\allowbreak\{\}COMPARISONS & 8460 & 0\\
games, 1000, Release date, ,  & avl & RECOLORS & 0 & 0\\
games, 1000, Release date, ,  & avl & ROTATIONS & 631 & 0\\
games, 1000, Release date, ,  & avl & TREE\_\allowbreak\{\}HEIGHT & 12 & 0\\
games, 1000, Genres, ,  & rubro\_\allowbreak\{\}negra & BUILD\_\allowbreak\{\}MS & 54.7364 & 26.6762\\
games, 1000, Genres, ,  & rubro\_\allowbreak\{\}negra & INSERT\_\allowbreak\{\}COMPARISONS & 6056 & 0\\
games, 1000, Genres, ,  & rubro\_\allowbreak\{\}negra & RECOLORS & 584 & 0\\
games, 1000, Genres, ,  & rubro\_\allowbreak\{\}negra & ROTATIONS & 149 & 0\\
games, 1000, Genres, ,  & rubro\_\allowbreak\{\}negra & TREE\_\allowbreak\{\}HEIGHT & 10 & 0\\
games, 1000, Release date, ,  & rubro\_\allowbreak\{\}negra & BUILD\_\allowbreak\{\}MS & 55.0765 & 27.5832\\
games, 1000, Release date, ,  & rubro\_\allowbreak\{\}negra & INSERT\_\allowbreak\{\}COMPARISONS & 8526 & 0\\
games, 1000, Release date, ,  & rubro\_\allowbreak\{\}negra & RECOLORS & 1919 & 0\\
games, 1000, Release date, ,  & rubro\_\allowbreak\{\}negra & ROTATIONS & 481 & 0\\
games, 1000, Release date, ,  & rubro\_\allowbreak\{\}negra & TREE\_\allowbreak\{\}HEIGHT & 12 & 0\\
games, 10000, Genres, ,  & abb & BUILD\_\allowbreak\{\}MS & 239.437 & 15.9696\\
games, 10000, Genres, ,  & abb & INSERT\_\allowbreak\{\}COMPARISONS & 84732 & 0\\
games, 10000, Genres, ,  & abb & RECOLORS & 0 & 0\\
games, 10000, Genres, ,  & abb & ROTATIONS & 0 & 0\\
games, 10000, Genres, ,  & abb & TREE\_\allowbreak\{\}HEIGHT & 24 & 0\\
games, 10000, Release date, ,  & abb & BUILD\_\allowbreak\{\}MS & 234.711 & 4.86402\\
games, 10000, Release date, ,  & abb & INSERT\_\allowbreak\{\}COMPARISONS & 132272 & 0\\
games, 10000, Release date, ,  & abb & RECOLORS & 0 & 0\\
games, 10000, Release date, ,  & abb & ROTATIONS & 0 & 0\\
games, 10000, Release date, ,  & abb & TREE\_\allowbreak\{\}HEIGHT & 28 & 0\\
games, 10000, Genres, ,  & avl & BUILD\_\allowbreak\{\}MS & 239.616 & 16.7397\\
games, 10000, Genres, ,  & avl & INSERT\_\allowbreak\{\}COMPARISONS & 77576 & 0\\
games, 10000, Genres, ,  & avl & RECOLORS & 0 & 0\\
games, 10000, Genres, ,  & avl & ROTATIONS & 627 & 0\\
games, 10000, Genres, ,  & avl & TREE\_\allowbreak\{\}HEIGHT & 11 & 0\\
games, 10000, Release date, ,  & avl & BUILD\_\allowbreak\{\}MS & 233.186 & 4.98049\\
games, 10000, Release date, ,  & avl & INSERT\_\allowbreak\{\}COMPARISONS & 103155 & 0\\
games, 10000, Release date, ,  & avl & RECOLORS & 0 & 0\\
games, 10000, Release date, ,  & avl & ROTATIONS & 2259 & 0\\
games, 10000, Release date, ,  & avl & TREE\_\allowbreak\{\}HEIGHT & 14 & 0\\
games, 10000, Genres, ,  & rubro\_\allowbreak\{\}negra & BUILD\_\allowbreak\{\}MS & 237.871 & 12.5673\\
games, 10000, Genres, ,  & rubro\_\allowbreak\{\}negra & INSERT\_\allowbreak\{\}COMPARISONS & 77020 & 0\\
games, 10000, Genres, ,  & rubro\_\allowbreak\{\}negra & RECOLORS & 2087 & 0\\
games, 10000, Genres, ,  & rubro\_\allowbreak\{\}negra & ROTATIONS & 527 & 0\\
games, 10000, Genres, ,  & rubro\_\allowbreak\{\}negra & TREE\_\allowbreak\{\}HEIGHT & 12 & 0\\
games, 10000, Release date, ,  & rubro\_\allowbreak\{\}negra & BUILD\_\allowbreak\{\}MS & 233.182 & 5.51643\\
games, 10000, Release date, ,  & rubro\_\allowbreak\{\}negra & INSERT\_\allowbreak\{\}COMPARISONS & 104699 & 0\\
games, 10000, Release date, ,  & rubro\_\allowbreak\{\}negra & RECOLORS & 7672 & 0\\
games, 10000, Release date, ,  & rubro\_\allowbreak\{\}negra & ROTATIONS & 1973 & 0\\
games, 10000, Release date, ,  & rubro\_\allowbreak\{\}negra & TREE\_\allowbreak\{\}HEIGHT & 14 & 0\\
\hline
\end{longtable}
\normalsize
}{\TODO{executar experimentos/analise.py para inserir a tabela de construção}}

\subsection{Busca simples e complexa}
\IfFileExists{tabela_busca.tex}{\scriptsize
\begin{longtable}{p{4cm}p{1.4cm}p{2.6cm}rr}
\hline
Base/configuração & Árvore & Métrica & Média & Desvio\\
\hline
\endfirsthead
\hline
Base/configuração & Árvore & Métrica & Média & Desvio\\
\hline
\endhead
asteroids, 1000, , complex, 10 & abb & FETCH\_\allowbreak\{\}MS & 2.4791 & 1.10915\\
asteroids, 1000, , complex, 10 & abb & LOAD\_\allowbreak\{\}INDEX\_\allowbreak\{\}MS & 0.739 & 0.282003\\
asteroids, 1000, , complex, 10 & abb & SEARCH\_\allowbreak\{\}COMPARISONS & 52 & 0\\
asteroids, 1000, , complex, 10 & abb & SEARCH\_\allowbreak\{\}MS & 0.0037 & 0.00149443\\
asteroids, 1000, , complex, 10 & abb & TOTAL\_\allowbreak\{\}MS & 3.9641 & 1.66841\\
asteroids, 1000, , complex, 10 & abb & WRITE\_\allowbreak\{\}MS & 0.3548 & 0.178459\\
asteroids, 1000, , complex, 10 & avl & FETCH\_\allowbreak\{\}MS & 2.7778 & 1.50281\\
asteroids, 1000, , complex, 10 & avl & LOAD\_\allowbreak\{\}INDEX\_\allowbreak\{\}MS & 0.807 & 0.436951\\
asteroids, 1000, , complex, 10 & avl & SEARCH\_\allowbreak\{\}COMPARISONS & 38 & 0\\
asteroids, 1000, , complex, 10 & avl & SEARCH\_\allowbreak\{\}MS & 0.0029 & 0.00119722\\
asteroids, 1000, , complex, 10 & avl & TOTAL\_\allowbreak\{\}MS & 4.3416 & 2.26868\\
asteroids, 1000, , complex, 10 & avl & WRITE\_\allowbreak\{\}MS & 0.3735 & 0.20602\\
asteroids, 1000, , complex, 10 & rubro\_\allowbreak\{\}negra & FETCH\_\allowbreak\{\}MS & 2.327 & 0.759055\\
asteroids, 1000, , complex, 10 & rubro\_\allowbreak\{\}negra & LOAD\_\allowbreak\{\}INDEX\_\allowbreak\{\}MS & 0.7516 & 0.212927\\
asteroids, 1000, , complex, 10 & rubro\_\allowbreak\{\}negra & SEARCH\_\allowbreak\{\}COMPARISONS & 37 & 0\\
asteroids, 1000, , complex, 10 & rubro\_\allowbreak\{\}negra & SEARCH\_\allowbreak\{\}MS & 0.0031 & 0.000875595\\
asteroids, 1000, , complex, 10 & rubro\_\allowbreak\{\}negra & TOTAL\_\allowbreak\{\}MS & 3.7757 & 1.17561\\
asteroids, 1000, , complex, 10 & rubro\_\allowbreak\{\}negra & WRITE\_\allowbreak\{\}MS & 0.3258 & 0.113623\\
asteroids, 1000, , complex, 100 & abb & FETCH\_\allowbreak\{\}MS & 2.7287 & 1.30192\\
asteroids, 1000, , complex, 100 & abb & LOAD\_\allowbreak\{\}INDEX\_\allowbreak\{\}MS & 0.7421 & 0.313017\\
asteroids, 1000, , complex, 100 & abb & SEARCH\_\allowbreak\{\}COMPARISONS & 388 & 0\\
asteroids, 1000, , complex, 100 & abb & SEARCH\_\allowbreak\{\}MS & 0.0207 & 0.00763108\\
asteroids, 1000, , complex, 100 & abb & TOTAL\_\allowbreak\{\}MS & 4.3922 & 2.00924\\
asteroids, 1000, , complex, 100 & abb & WRITE\_\allowbreak\{\}MS & 0.3899 & 0.204634\\
asteroids, 1000, , complex, 100 & avl & FETCH\_\allowbreak\{\}MS & 3.2321 & 1.66326\\
asteroids, 1000, , complex, 100 & avl & LOAD\_\allowbreak\{\}INDEX\_\allowbreak\{\}MS & 0.8688 & 0.407451\\
asteroids, 1000, , complex, 100 & avl & SEARCH\_\allowbreak\{\}COMPARISONS & 388 & 0\\
asteroids, 1000, , complex, 100 & avl & SEARCH\_\allowbreak\{\}MS & 0.0247 & 0.0119075\\
asteroids, 1000, , complex, 100 & avl & TOTAL\_\allowbreak\{\}MS & 5.2162 & 2.60915\\
asteroids, 1000, , complex, 100 & avl & WRITE\_\allowbreak\{\}MS & 0.5332 & 0.372776\\
asteroids, 1000, , complex, 100 & rubro\_\allowbreak\{\}negra & FETCH\_\allowbreak\{\}MS & 3.3274 & 1.79154\\
asteroids, 1000, , complex, 100 & rubro\_\allowbreak\{\}negra & LOAD\_\allowbreak\{\}INDEX\_\allowbreak\{\}MS & 0.8895 & 0.45373\\
asteroids, 1000, , complex, 100 & rubro\_\allowbreak\{\}negra & SEARCH\_\allowbreak\{\}COMPARISONS & 380 & 0\\
asteroids, 1000, , complex, 100 & rubro\_\allowbreak\{\}negra & SEARCH\_\allowbreak\{\}MS & 0.0208 & 0.00825698\\
asteroids, 1000, , complex, 100 & rubro\_\allowbreak\{\}negra & TOTAL\_\allowbreak\{\}MS & 5.2697 & 2.68408\\
asteroids, 1000, , complex, 100 & rubro\_\allowbreak\{\}negra & WRITE\_\allowbreak\{\}MS & 0.4649 & 0.253998\\
asteroids, 1000, , complex, 1000 & abb & FETCH\_\allowbreak\{\}MS & 2.9406 & 1.35218\\
asteroids, 1000, , complex, 1000 & abb & LOAD\_\allowbreak\{\}INDEX\_\allowbreak\{\}MS & 0.8991 & 0.391386\\
asteroids, 1000, , complex, 1000 & abb & SEARCH\_\allowbreak\{\}COMPARISONS & 4086 & 0\\
asteroids, 1000, , complex, 1000 & abb & SEARCH\_\allowbreak\{\}MS & 0.1958 & 0.0704222\\
asteroids, 1000, , complex, 1000 & abb & TOTAL\_\allowbreak\{\}MS & 5.9932 & 2.65214\\
asteroids, 1000, , complex, 1000 & abb & WRITE\_\allowbreak\{\}MS & 0.4387 & 0.247763\\
asteroids, 1000, , complex, 1000 & avl & FETCH\_\allowbreak\{\}MS & 2.9285 & 1.25755\\
asteroids, 1000, , complex, 1000 & avl & LOAD\_\allowbreak\{\}INDEX\_\allowbreak\{\}MS & 0.783 & 0.38138\\
asteroids, 1000, , complex, 1000 & avl & SEARCH\_\allowbreak\{\}COMPARISONS & 3598 & 0\\
asteroids, 1000, , complex, 1000 & avl & SEARCH\_\allowbreak\{\}MS & 0.1637 & 0.064474\\
asteroids, 1000, , complex, 1000 & avl & TOTAL\_\allowbreak\{\}MS & 5.7413 & 2.45343\\
asteroids, 1000, , complex, 1000 & avl & WRITE\_\allowbreak\{\}MS & 0.4761 & 0.227208\\
asteroids, 1000, , complex, 1000 & rubro\_\allowbreak\{\}negra & FETCH\_\allowbreak\{\}MS & 3.3412 & 1.5431\\
asteroids, 1000, , complex, 1000 & rubro\_\allowbreak\{\}negra & LOAD\_\allowbreak\{\}INDEX\_\allowbreak\{\}MS & 0.9515 & 0.42133\\
asteroids, 1000, , complex, 1000 & rubro\_\allowbreak\{\}negra & SEARCH\_\allowbreak\{\}COMPARISONS & 3620 & 0\\
asteroids, 1000, , complex, 1000 & rubro\_\allowbreak\{\}negra & SEARCH\_\allowbreak\{\}MS & 0.1842 & 0.0645442\\
asteroids, 1000, , complex, 1000 & rubro\_\allowbreak\{\}negra & TOTAL\_\allowbreak\{\}MS & 6.6461 & 2.9083\\
asteroids, 1000, , complex, 1000 & rubro\_\allowbreak\{\}negra & WRITE\_\allowbreak\{\}MS & 0.4628 & 0.213039\\
asteroids, 1000, , simple, 1 & abb & FETCH\_\allowbreak\{\}MS & 0.1574 & 0.0622293\\
asteroids, 1000, , simple, 1 & abb & LOAD\_\allowbreak\{\}INDEX\_\allowbreak\{\}MS & 0.527 & 0.221365\\
asteroids, 1000, , simple, 1 & abb & SEARCH\_\allowbreak\{\}COMPARISONS & 5 & 0\\
asteroids, 1000, , simple, 1 & abb & SEARCH\_\allowbreak\{\}MS & 0.001 & 0.000471405\\
asteroids, 1000, , simple, 1 & abb & TOTAL\_\allowbreak\{\}MS & 0.8212 & 0.34177\\
asteroids, 1000, , simple, 1 & abb & WRITE\_\allowbreak\{\}MS & 0.0415 & 0.0183742\\
asteroids, 1000, , simple, 1 & avl & FETCH\_\allowbreak\{\}MS & 0.1873 & 0.0964792\\
asteroids, 1000, , simple, 1 & avl & LOAD\_\allowbreak\{\}INDEX\_\allowbreak\{\}MS & 0.5659 & 0.292987\\
asteroids, 1000, , simple, 1 & avl & SEARCH\_\allowbreak\{\}COMPARISONS & 6 & 0\\
asteroids, 1000, , simple, 1 & avl & SEARCH\_\allowbreak\{\}MS & 0.0007 & 0.000483046\\
asteroids, 1000, , simple, 1 & avl & TOTAL\_\allowbreak\{\}MS & 0.8976 & 0.450204\\
asteroids, 1000, , simple, 1 & avl & WRITE\_\allowbreak\{\}MS & 0.0471 & 0.0211894\\
asteroids, 1000, , simple, 1 & rubro\_\allowbreak\{\}negra & FETCH\_\allowbreak\{\}MS & 0.2238 & 0.140298\\
asteroids, 1000, , simple, 1 & rubro\_\allowbreak\{\}negra & LOAD\_\allowbreak\{\}INDEX\_\allowbreak\{\}MS & 0.6408 & 0.321387\\
asteroids, 1000, , simple, 1 & rubro\_\allowbreak\{\}negra & SEARCH\_\allowbreak\{\}COMPARISONS & 7 & 0\\
asteroids, 1000, , simple, 1 & rubro\_\allowbreak\{\}negra & SEARCH\_\allowbreak\{\}MS & 0.0008 & 0.000632456\\
asteroids, 1000, , simple, 1 & rubro\_\allowbreak\{\}negra & TOTAL\_\allowbreak\{\}MS & 1.0307 & 0.526885\\
asteroids, 1000, , simple, 1 & rubro\_\allowbreak\{\}negra & WRITE\_\allowbreak\{\}MS & 0.0532 & 0.0244259\\
asteroids, 1000, , simple, 10 & abb & FETCH\_\allowbreak\{\}MS & 0.3125 & 0.133765\\
asteroids, 1000, , simple, 10 & abb & LOAD\_\allowbreak\{\}INDEX\_\allowbreak\{\}MS & 0.5628 & 0.215041\\
asteroids, 1000, , simple, 10 & abb & SEARCH\_\allowbreak\{\}COMPARISONS & 63 & 0\\
asteroids, 1000, , simple, 10 & abb & SEARCH\_\allowbreak\{\}MS & 0.0035 & 0.0012693\\
asteroids, 1000, , simple, 10 & abb & TOTAL\_\allowbreak\{\}MS & 1.0804 & 0.437137\\
asteroids, 1000, , simple, 10 & abb & WRITE\_\allowbreak\{\}MS & 0.0724 & 0.0419714\\
asteroids, 1000, , simple, 10 & avl & FETCH\_\allowbreak\{\}MS & 0.2889 & 0.142628\\
asteroids, 1000, , simple, 10 & avl & LOAD\_\allowbreak\{\}INDEX\_\allowbreak\{\}MS & 0.5353 & 0.24412\\
asteroids, 1000, , simple, 10 & avl & SEARCH\_\allowbreak\{\}COMPARISONS & 65 & 0\\
asteroids, 1000, , simple, 10 & avl & SEARCH\_\allowbreak\{\}MS & 0.0038 & 0.00154919\\
asteroids, 1000, , simple, 10 & avl & TOTAL\_\allowbreak\{\}MS & 1.0063 & 0.470629\\
asteroids, 1000, , simple, 10 & avl & WRITE\_\allowbreak\{\}MS & 0.06 & 0.031756\\
asteroids, 1000, , simple, 10 & rubro\_\allowbreak\{\}negra & FETCH\_\allowbreak\{\}MS & 0.293 & 0.123476\\
asteroids, 1000, , simple, 10 & rubro\_\allowbreak\{\}negra & LOAD\_\allowbreak\{\}INDEX\_\allowbreak\{\}MS & 0.5334 & 0.2311\\
asteroids, 1000, , simple, 10 & rubro\_\allowbreak\{\}negra & SEARCH\_\allowbreak\{\}COMPARISONS & 58 & 0\\
asteroids, 1000, , simple, 10 & rubro\_\allowbreak\{\}negra & SEARCH\_\allowbreak\{\}MS & 0.0034 & 0.00142984\\
asteroids, 1000, , simple, 10 & rubro\_\allowbreak\{\}negra & TOTAL\_\allowbreak\{\}MS & 1.0165 & 0.441362\\
asteroids, 1000, , simple, 10 & rubro\_\allowbreak\{\}negra & WRITE\_\allowbreak\{\}MS & 0.0585 & 0.0300046\\
asteroids, 1000, , simple, 100 & abb & FETCH\_\allowbreak\{\}MS & 1.331 & 0.694868\\
asteroids, 1000, , simple, 100 & abb & LOAD\_\allowbreak\{\}INDEX\_\allowbreak\{\}MS & 0.5328 & 0.25711\\
asteroids, 1000, , simple, 100 & abb & SEARCH\_\allowbreak\{\}COMPARISONS & 667 & 0\\
asteroids, 1000, , simple, 100 & abb & SEARCH\_\allowbreak\{\}MS & 0.0299 & 0.012078\\
asteroids, 1000, , simple, 100 & abb & TOTAL\_\allowbreak\{\}MS & 2.4496 & 1.21842\\
asteroids, 1000, , simple, 100 & abb & WRITE\_\allowbreak\{\}MS & 0.2147 & 0.122122\\
asteroids, 1000, , simple, 100 & avl & FETCH\_\allowbreak\{\}MS & 1.327 & 0.67856\\
asteroids, 1000, , simple, 100 & avl & LOAD\_\allowbreak\{\}INDEX\_\allowbreak\{\}MS & 0.6191 & 0.296175\\
asteroids, 1000, , simple, 100 & avl & SEARCH\_\allowbreak\{\}COMPARISONS & 652 & 0\\
asteroids, 1000, , simple, 100 & avl & SEARCH\_\allowbreak\{\}MS & 0.0283 & 0.0120282\\
asteroids, 1000, , simple, 100 & avl & TOTAL\_\allowbreak\{\}MS & 2.5333 & 1.23307\\
asteroids, 1000, , simple, 100 & avl & WRITE\_\allowbreak\{\}MS & 0.1965 & 0.100125\\
asteroids, 1000, , simple, 100 & rubro\_\allowbreak\{\}negra & FETCH\_\allowbreak\{\}MS & 1.3841 & 0.817396\\
asteroids, 1000, , simple, 100 & rubro\_\allowbreak\{\}negra & LOAD\_\allowbreak\{\}INDEX\_\allowbreak\{\}MS & 0.5986 & 0.293385\\
asteroids, 1000, , simple, 100 & rubro\_\allowbreak\{\}negra & SEARCH\_\allowbreak\{\}COMPARISONS & 632 & 0\\
asteroids, 1000, , simple, 100 & rubro\_\allowbreak\{\}negra & SEARCH\_\allowbreak\{\}MS & 0.0274 & 0.0114232\\
asteroids, 1000, , simple, 100 & rubro\_\allowbreak\{\}negra & TOTAL\_\allowbreak\{\}MS & 2.5824 & 1.41944\\
asteroids, 1000, , simple, 100 & rubro\_\allowbreak\{\}negra & WRITE\_\allowbreak\{\}MS & 0.2116 & 0.122247\\
asteroids, 1000, , simple, 1000 & abb & FETCH\_\allowbreak\{\}MS & 2.8433 & 1.36317\\
asteroids, 1000, , simple, 1000 & abb & LOAD\_\allowbreak\{\}INDEX\_\allowbreak\{\}MS & 0.543 & 0.188658\\
asteroids, 1000, , simple, 1000 & abb & SEARCH\_\allowbreak\{\}COMPARISONS & 6530 & 0\\
asteroids, 1000, , simple, 1000 & abb & SEARCH\_\allowbreak\{\}MS & 0.2723 & 0.0979161\\
asteroids, 1000, , simple, 1000 & abb & TOTAL\_\allowbreak\{\}MS & 5.3445 & 2.33721\\
asteroids, 1000, , simple, 1000 & abb & WRITE\_\allowbreak\{\}MS & 0.4146 & 0.222464\\
asteroids, 1000, , simple, 1000 & avl & FETCH\_\allowbreak\{\}MS & 3.2277 & 1.47156\\
asteroids, 1000, , simple, 1000 & avl & LOAD\_\allowbreak\{\}INDEX\_\allowbreak\{\}MS & 0.65 & 0.284092\\
asteroids, 1000, , simple, 1000 & avl & SEARCH\_\allowbreak\{\}COMPARISONS & 6278 & 0\\
asteroids, 1000, , simple, 1000 & avl & SEARCH\_\allowbreak\{\}MS & 0.2457 & 0.0885162\\
asteroids, 1000, , simple, 1000 & avl & TOTAL\_\allowbreak\{\}MS & 6.2496 & 3.10581\\
asteroids, 1000, , simple, 1000 & avl & WRITE\_\allowbreak\{\}MS & 0.4755 & 0.253369\\
asteroids, 1000, , simple, 1000 & rubro\_\allowbreak\{\}negra & FETCH\_\allowbreak\{\}MS & 2.7439 & 1.40695\\
asteroids, 1000, , simple, 1000 & rubro\_\allowbreak\{\}negra & LOAD\_\allowbreak\{\}INDEX\_\allowbreak\{\}MS & 0.6099 & 0.320833\\
asteroids, 1000, , simple, 1000 & rubro\_\allowbreak\{\}negra & SEARCH\_\allowbreak\{\}COMPARISONS & 5964 & 0\\
asteroids, 1000, , simple, 1000 & rubro\_\allowbreak\{\}negra & SEARCH\_\allowbreak\{\}MS & 0.2292 & 0.0873496\\
asteroids, 1000, , simple, 1000 & rubro\_\allowbreak\{\}negra & TOTAL\_\allowbreak\{\}MS & 5.2535 & 2.59922\\
asteroids, 1000, , simple, 1000 & rubro\_\allowbreak\{\}negra & WRITE\_\allowbreak\{\}MS & 0.3693 & 0.23179\\
asteroids, 10000, , complex, 10 & abb & FETCH\_\allowbreak\{\}MS & 34.7797 & 3.64245\\
asteroids, 10000, , complex, 10 & abb & LOAD\_\allowbreak\{\}INDEX\_\allowbreak\{\}MS & 3.3894 & 0.962822\\
asteroids, 10000, , complex, 10 & abb & SEARCH\_\allowbreak\{\}COMPARISONS & 59 & 0\\
asteroids, 10000, , complex, 10 & abb & SEARCH\_\allowbreak\{\}MS & 0.0059 & 0.000994429\\
asteroids, 10000, , complex, 10 & abb & TOTAL\_\allowbreak\{\}MS & 46.4373 & 4.49728\\
asteroids, 10000, , complex, 10 & abb & WRITE\_\allowbreak\{\}MS & 4.6684 & 0.602082\\
asteroids, 10000, , complex, 10 & avl & FETCH\_\allowbreak\{\}MS & 35.519 & 5.17217\\
asteroids, 10000, , complex, 10 & avl & LOAD\_\allowbreak\{\}INDEX\_\allowbreak\{\}MS & 3.3078 & 0.780627\\
asteroids, 10000, , complex, 10 & avl & SEARCH\_\allowbreak\{\}COMPARISONS & 39 & 0\\
asteroids, 10000, , complex, 10 & avl & SEARCH\_\allowbreak\{\}MS & 0.0055 & 0.000849837\\
asteroids, 10000, , complex, 10 & avl & TOTAL\_\allowbreak\{\}MS & 47.3492 & 6.17105\\
asteroids, 10000, , complex, 10 & avl & WRITE\_\allowbreak\{\}MS & 4.795 & 0.735116\\
asteroids, 10000, , complex, 10 & rubro\_\allowbreak\{\}negra & FETCH\_\allowbreak\{\}MS & 32.9128 & 4.73131\\
asteroids, 10000, , complex, 10 & rubro\_\allowbreak\{\}negra & LOAD\_\allowbreak\{\}INDEX\_\allowbreak\{\}MS & 3.1094 & 0.770999\\
asteroids, 10000, , complex, 10 & rubro\_\allowbreak\{\}negra & SEARCH\_\allowbreak\{\}COMPARISONS & 39 & 0\\
asteroids, 10000, , complex, 10 & rubro\_\allowbreak\{\}negra & SEARCH\_\allowbreak\{\}MS & 0.0045 & 0.000527046\\
asteroids, 10000, , complex, 10 & rubro\_\allowbreak\{\}negra & TOTAL\_\allowbreak\{\}MS & 44.0307 & 6.00862\\
asteroids, 10000, , complex, 10 & rubro\_\allowbreak\{\}negra & WRITE\_\allowbreak\{\}MS & 4.4009 & 0.753899\\
asteroids, 10000, , complex, 100 & abb & FETCH\_\allowbreak\{\}MS & 37.222 & 4.24496\\
asteroids, 10000, , complex, 100 & abb & LOAD\_\allowbreak\{\}INDEX\_\allowbreak\{\}MS & 3.2592 & 0.794307\\
asteroids, 10000, , complex, 100 & abb & SEARCH\_\allowbreak\{\}COMPARISONS & 434 & 0\\
asteroids, 10000, , complex, 100 & abb & SEARCH\_\allowbreak\{\}MS & 0.0328 & 0.0039101\\
asteroids, 10000, , complex, 100 & abb & TOTAL\_\allowbreak\{\}MS & 50.2983 & 5.68577\\
asteroids, 10000, , complex, 100 & abb & WRITE\_\allowbreak\{\}MS & 5.0076 & 0.582162\\
asteroids, 10000, , complex, 100 & avl & FETCH\_\allowbreak\{\}MS & 36.5476 & 4.34737\\
asteroids, 10000, , complex, 100 & avl & LOAD\_\allowbreak\{\}INDEX\_\allowbreak\{\}MS & 3.1972 & 1.52588\\
asteroids, 10000, , complex, 100 & avl & SEARCH\_\allowbreak\{\}COMPARISONS & 358 & 0\\
asteroids, 10000, , complex, 100 & avl & SEARCH\_\allowbreak\{\}MS & 0.0281 & 0.0048637\\
asteroids, 10000, , complex, 100 & avl & TOTAL\_\allowbreak\{\}MS & 49.1999 & 4.34161\\
asteroids, 10000, , complex, 100 & avl & WRITE\_\allowbreak\{\}MS & 4.8414 & 0.658185\\
asteroids, 10000, , complex, 100 & rubro\_\allowbreak\{\}negra & FETCH\_\allowbreak\{\}MS & 32.3353 & 8.43442\\
asteroids, 10000, , complex, 100 & rubro\_\allowbreak\{\}negra & LOAD\_\allowbreak\{\}INDEX\_\allowbreak\{\}MS & 3.4234 & 0.952096\\
asteroids, 10000, , complex, 100 & rubro\_\allowbreak\{\}negra & SEARCH\_\allowbreak\{\}COMPARISONS & 356 & 0\\
asteroids, 10000, , complex, 100 & rubro\_\allowbreak\{\}negra & SEARCH\_\allowbreak\{\}MS & 0.0282 & 0.00547317\\
asteroids, 10000, , complex, 100 & rubro\_\allowbreak\{\}negra & TOTAL\_\allowbreak\{\}MS & 44.5258 & 10.6265\\
asteroids, 10000, , complex, 100 & rubro\_\allowbreak\{\}negra & WRITE\_\allowbreak\{\}MS & 4.2696 & 1.20098\\
asteroids, 10000, , complex, 1000 & abb & FETCH\_\allowbreak\{\}MS & 33.4539 & 9.39035\\
asteroids, 10000, , complex, 1000 & abb & LOAD\_\allowbreak\{\}INDEX\_\allowbreak\{\}MS & 2.8945 & 0.858659\\
asteroids, 10000, , complex, 1000 & abb & SEARCH\_\allowbreak\{\}COMPARISONS & 4528 & 0\\
asteroids, 10000, , complex, 1000 & abb & SEARCH\_\allowbreak\{\}MS & 0.2585 & 0.0577143\\
asteroids, 10000, , complex, 1000 & abb & TOTAL\_\allowbreak\{\}MS & 48.0052 & 12.2145\\
asteroids, 10000, , complex, 1000 & abb & WRITE\_\allowbreak\{\}MS & 4.418 & 1.29849\\
asteroids, 10000, , complex, 1000 & avl & FETCH\_\allowbreak\{\}MS & 34.2614 & 7.54924\\
asteroids, 10000, , complex, 1000 & avl & LOAD\_\allowbreak\{\}INDEX\_\allowbreak\{\}MS & 3.3045 & 1.23402\\
asteroids, 10000, , complex, 1000 & avl & SEARCH\_\allowbreak\{\}COMPARISONS & 3684 & 0\\
asteroids, 10000, , complex, 1000 & avl & SEARCH\_\allowbreak\{\}MS & 0.246 & 0.0878964\\
asteroids, 10000, , complex, 1000 & avl & TOTAL\_\allowbreak\{\}MS & 49.8336 & 10.8923\\
asteroids, 10000, , complex, 1000 & avl & WRITE\_\allowbreak\{\}MS & 4.6517 & 1.07563\\
asteroids, 10000, , complex, 1000 & rubro\_\allowbreak\{\}negra & FETCH\_\allowbreak\{\}MS & 34.8828 & 4.7432\\
asteroids, 10000, , complex, 1000 & rubro\_\allowbreak\{\}negra & LOAD\_\allowbreak\{\}INDEX\_\allowbreak\{\}MS & 3.1912 & 0.982891\\
asteroids, 10000, , complex, 1000 & rubro\_\allowbreak\{\}negra & SEARCH\_\allowbreak\{\}COMPARISONS & 3720 & 0\\
asteroids, 10000, , complex, 1000 & rubro\_\allowbreak\{\}negra & SEARCH\_\allowbreak\{\}MS & 0.2548 & 0.0676294\\
asteroids, 10000, , complex, 1000 & rubro\_\allowbreak\{\}negra & TOTAL\_\allowbreak\{\}MS & 50.4643 & 6.6934\\
asteroids, 10000, , complex, 1000 & rubro\_\allowbreak\{\}negra & WRITE\_\allowbreak\{\}MS & 4.6217 & 0.744667\\
asteroids, 10000, , simple, 1 & abb & FETCH\_\allowbreak\{\}MS & 1.8623 & 0.470492\\
asteroids, 10000, , simple, 1 & abb & LOAD\_\allowbreak\{\}INDEX\_\allowbreak\{\}MS & 1.5162 & 0.409164\\
asteroids, 10000, , simple, 1 & abb & SEARCH\_\allowbreak\{\}COMPARISONS & 5 & 0\\
asteroids, 10000, , simple, 1 & abb & SEARCH\_\allowbreak\{\}MS & 0.0013 & 0.000483046\\
asteroids, 10000, , simple, 1 & abb & TOTAL\_\allowbreak\{\}MS & 3.8293 & 0.962279\\
asteroids, 10000, , simple, 1 & abb & WRITE\_\allowbreak\{\}MS & 0.2131 & 0.0540606\\
asteroids, 10000, , simple, 1 & avl & FETCH\_\allowbreak\{\}MS & 1.919 & 0.631065\\
asteroids, 10000, , simple, 1 & avl & LOAD\_\allowbreak\{\}INDEX\_\allowbreak\{\}MS & 1.8021 & 0.58863\\
asteroids, 10000, , simple, 1 & avl & SEARCH\_\allowbreak\{\}COMPARISONS & 7 & 0\\
asteroids, 10000, , simple, 1 & avl & SEARCH\_\allowbreak\{\}MS & 0.0014 & 0.000516398\\
asteroids, 10000, , simple, 1 & avl & TOTAL\_\allowbreak\{\}MS & 4.2016 & 1.3383\\
asteroids, 10000, , simple, 1 & avl & WRITE\_\allowbreak\{\}MS & 0.2211 & 0.0840508\\
asteroids, 10000, , simple, 1 & rubro\_\allowbreak\{\}negra & FETCH\_\allowbreak\{\}MS & 1.8953 & 0.603277\\
asteroids, 10000, , simple, 1 & rubro\_\allowbreak\{\}negra & LOAD\_\allowbreak\{\}INDEX\_\allowbreak\{\}MS & 1.9155 & 1.17918\\
asteroids, 10000, , simple, 1 & rubro\_\allowbreak\{\}negra & SEARCH\_\allowbreak\{\}COMPARISONS & 6 & 0\\
asteroids, 10000, , simple, 1 & rubro\_\allowbreak\{\}negra & SEARCH\_\allowbreak\{\}MS & 0.0013 & 0.000674949\\
asteroids, 10000, , simple, 1 & rubro\_\allowbreak\{\}negra & TOTAL\_\allowbreak\{\}MS & 4.2846 & 1.83347\\
asteroids, 10000, , simple, 1 & rubro\_\allowbreak\{\}negra & WRITE\_\allowbreak\{\}MS & 0.2153 & 0.0720155\\
asteroids, 10000, , simple, 10 & abb & FETCH\_\allowbreak\{\}MS & 3.4439 & 1.00169\\
asteroids, 10000, , simple, 10 & abb & LOAD\_\allowbreak\{\}INDEX\_\allowbreak\{\}MS & 1.7335 & 0.615747\\
asteroids, 10000, , simple, 10 & abb & SEARCH\_\allowbreak\{\}COMPARISONS & 77 & 0\\
asteroids, 10000, , simple, 10 & abb & SEARCH\_\allowbreak\{\}MS & 0.0061 & 0.00128668\\
asteroids, 10000, , simple, 10 & abb & TOTAL\_\allowbreak\{\}MS & 6.0163 & 1.75021\\
asteroids, 10000, , simple, 10 & abb & WRITE\_\allowbreak\{\}MS & 0.4208 & 0.16734\\
asteroids, 10000, , simple, 10 & avl & FETCH\_\allowbreak\{\}MS & 3.6021 & 0.984024\\
asteroids, 10000, , simple, 10 & avl & LOAD\_\allowbreak\{\}INDEX\_\allowbreak\{\}MS & 1.7051 & 0.583326\\
asteroids, 10000, , simple, 10 & avl & SEARCH\_\allowbreak\{\}COMPARISONS & 56 & 0\\
asteroids, 10000, , simple, 10 & avl & SEARCH\_\allowbreak\{\}MS & 0.0051 & 0.00144914\\
asteroids, 10000, , simple, 10 & avl & TOTAL\_\allowbreak\{\}MS & 6.1947 & 1.7274\\
asteroids, 10000, , simple, 10 & avl & WRITE\_\allowbreak\{\}MS & 0.4556 & 0.146986\\
asteroids, 10000, , simple, 10 & rubro\_\allowbreak\{\}negra & FETCH\_\allowbreak\{\}MS & 3.2636 & 0.939203\\
asteroids, 10000, , simple, 10 & rubro\_\allowbreak\{\}negra & LOAD\_\allowbreak\{\}INDEX\_\allowbreak\{\}MS & 1.6492 & 0.567332\\
asteroids, 10000, , simple, 10 & rubro\_\allowbreak\{\}negra & SEARCH\_\allowbreak\{\}COMPARISONS & 51 & 0\\
asteroids, 10000, , simple, 10 & rubro\_\allowbreak\{\}negra & SEARCH\_\allowbreak\{\}MS & 0.0042 & 0.000918937\\
asteroids, 10000, , simple, 10 & rubro\_\allowbreak\{\}negra & TOTAL\_\allowbreak\{\}MS & 5.8182 & 1.7416\\
asteroids, 10000, , simple, 10 & rubro\_\allowbreak\{\}negra & WRITE\_\allowbreak\{\}MS & 0.4892 & 0.350095\\
asteroids, 10000, , simple, 100 & abb & FETCH\_\allowbreak\{\}MS & 16.2381 & 4.45898\\
asteroids, 10000, , simple, 100 & abb & LOAD\_\allowbreak\{\}INDEX\_\allowbreak\{\}MS & 1.6618 & 0.583277\\
asteroids, 10000, , simple, 100 & abb & SEARCH\_\allowbreak\{\}COMPARISONS & 743 & 0\\
asteroids, 10000, , simple, 100 & abb & SEARCH\_\allowbreak\{\}MS & 0.0407 & 0.00936957\\
asteroids, 10000, , simple, 100 & abb & TOTAL\_\allowbreak\{\}MS & 22.2525 & 5.92184\\
asteroids, 10000, , simple, 100 & abb & WRITE\_\allowbreak\{\}MS & 2.1129 & 0.598712\\
asteroids, 10000, , simple, 100 & avl & FETCH\_\allowbreak\{\}MS & 17.9321 & 4.16511\\
asteroids, 10000, , simple, 100 & avl & LOAD\_\allowbreak\{\}INDEX\_\allowbreak\{\}MS & 1.8423 & 0.501197\\
asteroids, 10000, , simple, 100 & avl & SEARCH\_\allowbreak\{\}COMPARISONS & 575 & 0\\
asteroids, 10000, , simple, 100 & avl & SEARCH\_\allowbreak\{\}MS & 0.0351 & 0.00462961\\
asteroids, 10000, , simple, 100 & avl & TOTAL\_\allowbreak\{\}MS & 24.3296 & 5.40169\\
asteroids, 10000, , simple, 100 & avl & WRITE\_\allowbreak\{\}MS & 2.2979 & 0.562598\\
asteroids, 10000, , simple, 100 & rubro\_\allowbreak\{\}negra & FETCH\_\allowbreak\{\}MS & 18.3336 & 3.69954\\
asteroids, 10000, , simple, 100 & rubro\_\allowbreak\{\}negra & LOAD\_\allowbreak\{\}INDEX\_\allowbreak\{\}MS & 1.8151 & 0.468557\\
asteroids, 10000, , simple, 100 & rubro\_\allowbreak\{\}negra & SEARCH\_\allowbreak\{\}COMPARISONS & 588 & 0\\
asteroids, 10000, , simple, 100 & rubro\_\allowbreak\{\}negra & SEARCH\_\allowbreak\{\}MS & 0.0363 & 0.00371334\\
asteroids, 10000, , simple, 100 & rubro\_\allowbreak\{\}negra & TOTAL\_\allowbreak\{\}MS & 25.0068 & 4.74507\\
asteroids, 10000, , simple, 100 & rubro\_\allowbreak\{\}negra & WRITE\_\allowbreak\{\}MS & 2.4804 & 0.618633\\
asteroids, 10000, , simple, 1000 & abb & FETCH\_\allowbreak\{\}MS & 36.8035 & 5.10058\\
asteroids, 10000, , simple, 1000 & abb & LOAD\_\allowbreak\{\}INDEX\_\allowbreak\{\}MS & 2.2522 & 0.555837\\
asteroids, 10000, , simple, 1000 & abb & SEARCH\_\allowbreak\{\}COMPARISONS & 7114 & 0\\
asteroids, 10000, , simple, 1000 & abb & SEARCH\_\allowbreak\{\}MS & 0.4028 & 0.051779\\
asteroids, 10000, , simple, 1000 & abb & TOTAL\_\allowbreak\{\}MS & 50.7863 & 5.29797\\
asteroids, 10000, , simple, 1000 & abb & WRITE\_\allowbreak\{\}MS & 4.8605 & 0.78664\\
asteroids, 10000, , simple, 1000 & avl & FETCH\_\allowbreak\{\}MS & 35.3246 & 3.76582\\
asteroids, 10000, , simple, 1000 & avl & LOAD\_\allowbreak\{\}INDEX\_\allowbreak\{\}MS & 1.7789 & 0.474454\\
asteroids, 10000, , simple, 1000 & avl & SEARCH\_\allowbreak\{\}COMPARISONS & 5602 & 0\\
asteroids, 10000, , simple, 1000 & avl & SEARCH\_\allowbreak\{\}MS & 0.2975 & 0.0471693\\
asteroids, 10000, , simple, 1000 & avl & TOTAL\_\allowbreak\{\}MS & 48.2095 & 4.14874\\
asteroids, 10000, , simple, 1000 & avl & WRITE\_\allowbreak\{\}MS & 4.7015 & 0.625269\\
asteroids, 10000, , simple, 1000 & rubro\_\allowbreak\{\}negra & FETCH\_\allowbreak\{\}MS & 36.5102 & 5.89517\\
asteroids, 10000, , simple, 1000 & rubro\_\allowbreak\{\}negra & LOAD\_\allowbreak\{\}INDEX\_\allowbreak\{\}MS & 1.69 & 0.449731\\
asteroids, 10000, , simple, 1000 & rubro\_\allowbreak\{\}negra & SEARCH\_\allowbreak\{\}COMPARISONS & 5656 & 0\\
asteroids, 10000, , simple, 1000 & rubro\_\allowbreak\{\}negra & SEARCH\_\allowbreak\{\}MS & 0.2856 & 0.0271015\\
asteroids, 10000, , simple, 1000 & rubro\_\allowbreak\{\}negra & TOTAL\_\allowbreak\{\}MS & 49.1559 & 6.78104\\
asteroids, 10000, , simple, 1000 & rubro\_\allowbreak\{\}negra & WRITE\_\allowbreak\{\}MS & 4.8495 & 0.859297\\
games, 1000, , complex, 10 & abb & FETCH\_\allowbreak\{\}MS & 0.4647 & 0.288697\\
games, 1000, , complex, 10 & abb & LOAD\_\allowbreak\{\}INDEX\_\allowbreak\{\}MS & 5.1584 & 2.60897\\
games, 1000, , complex, 10 & abb & SEARCH\_\allowbreak\{\}COMPARISONS & 83 & 0\\
games, 1000, , complex, 10 & abb & SEARCH\_\allowbreak\{\}MS & 0.0059 & 0.00284605\\
games, 1000, , complex, 10 & abb & TOTAL\_\allowbreak\{\}MS & 6.097 & 3.14247\\
games, 1000, , complex, 10 & abb & WRITE\_\allowbreak\{\}MS & 0.0983 & 0.102593\\
games, 1000, , complex, 10 & avl & FETCH\_\allowbreak\{\}MS & 0.4227 & 0.28021\\
games, 1000, , complex, 10 & avl & LOAD\_\allowbreak\{\}INDEX\_\allowbreak\{\}MS & 5.3494 & 3.31254\\
games, 1000, , complex, 10 & avl & SEARCH\_\allowbreak\{\}COMPARISONS & 88 & 0\\
games, 1000, , complex, 10 & avl & SEARCH\_\allowbreak\{\}MS & 0.0063 & 0.00352924\\
games, 1000, , complex, 10 & avl & TOTAL\_\allowbreak\{\}MS & 6.2239 & 3.86235\\
games, 1000, , complex, 10 & avl & WRITE\_\allowbreak\{\}MS & 0.0729 & 0.0462564\\
games, 1000, , complex, 10 & rubro\_\allowbreak\{\}negra & FETCH\_\allowbreak\{\}MS & 0.4476 & 0.247903\\
games, 1000, , complex, 10 & rubro\_\allowbreak\{\}negra & LOAD\_\allowbreak\{\}INDEX\_\allowbreak\{\}MS & 5.3551 & 2.97042\\
games, 1000, , complex, 10 & rubro\_\allowbreak\{\}negra & SEARCH\_\allowbreak\{\}COMPARISONS & 84 & 0\\
games, 1000, , complex, 10 & rubro\_\allowbreak\{\}negra & SEARCH\_\allowbreak\{\}MS & 0.0058 & 0.00248551\\
games, 1000, , complex, 10 & rubro\_\allowbreak\{\}negra & TOTAL\_\allowbreak\{\}MS & 6.2712 & 3.43996\\
games, 1000, , complex, 10 & rubro\_\allowbreak\{\}negra & WRITE\_\allowbreak\{\}MS & 0.081 & 0.0412634\\
games, 1000, , complex, 100 & abb & FETCH\_\allowbreak\{\}MS & 3.6751 & 2.2849\\
games, 1000, , complex, 100 & abb & LOAD\_\allowbreak\{\}INDEX\_\allowbreak\{\}MS & 4.9232 & 2.8401\\
games, 1000, , complex, 100 & abb & SEARCH\_\allowbreak\{\}COMPARISONS & 966 & 0\\
games, 1000, , complex, 100 & abb & SEARCH\_\allowbreak\{\}MS & 0.0434 & 0.0205058\\
games, 1000, , complex, 100 & abb & TOTAL\_\allowbreak\{\}MS & 9.4549 & 5.47358\\
games, 1000, , complex, 100 & abb & WRITE\_\allowbreak\{\}MS & 0.4055 & 0.255381\\
games, 1000, , complex, 100 & avl & FETCH\_\allowbreak\{\}MS & 3.819 & 2.06331\\
games, 1000, , complex, 100 & avl & LOAD\_\allowbreak\{\}INDEX\_\allowbreak\{\}MS & 4.9634 & 2.64521\\
games, 1000, , complex, 100 & avl & SEARCH\_\allowbreak\{\}COMPARISONS & 820 & 0\\
games, 1000, , complex, 100 & avl & SEARCH\_\allowbreak\{\}MS & 0.0382 & 0.0165784\\
games, 1000, , complex, 100 & avl & TOTAL\_\allowbreak\{\}MS & 9.7002 & 5.05646\\
games, 1000, , complex, 100 & avl & WRITE\_\allowbreak\{\}MS & 0.4256 & 0.232079\\
games, 1000, , complex, 100 & rubro\_\allowbreak\{\}negra & FETCH\_\allowbreak\{\}MS & 4.0316 & 2.37362\\
games, 1000, , complex, 100 & rubro\_\allowbreak\{\}negra & LOAD\_\allowbreak\{\}INDEX\_\allowbreak\{\}MS & 5.182 & 2.46832\\
games, 1000, , complex, 100 & rubro\_\allowbreak\{\}negra & SEARCH\_\allowbreak\{\}COMPARISONS & 810 & 0\\
games, 1000, , complex, 100 & rubro\_\allowbreak\{\}negra & SEARCH\_\allowbreak\{\}MS & 0.039 & 0.0161933\\
games, 1000, , complex, 100 & rubro\_\allowbreak\{\}negra & TOTAL\_\allowbreak\{\}MS & 10.1933 & 5.33938\\
games, 1000, , complex, 100 & rubro\_\allowbreak\{\}negra & WRITE\_\allowbreak\{\}MS & 0.4591 & 0.295561\\
games, 1000, , complex, 1000 & abb & FETCH\_\allowbreak\{\}MS & 23.5595 & 11.8625\\
games, 1000, , complex, 1000 & abb & LOAD\_\allowbreak\{\}INDEX\_\allowbreak\{\}MS & 5.8603 & 3.46427\\
games, 1000, , complex, 1000 & abb & SEARCH\_\allowbreak\{\}COMPARISONS & 9999 & 0\\
games, 1000, , complex, 1000 & abb & SEARCH\_\allowbreak\{\}MS & 0.4261 & 0.210699\\
games, 1000, , complex, 1000 & abb & TOTAL\_\allowbreak\{\}MS & 34.0306 & 17.6726\\
games, 1000, , complex, 1000 & abb & WRITE\_\allowbreak\{\}MS & 2.6333 & 1.44139\\
games, 1000, , complex, 1000 & avl & FETCH\_\allowbreak\{\}MS & 21.2256 & 10.5265\\
games, 1000, , complex, 1000 & avl & LOAD\_\allowbreak\{\}INDEX\_\allowbreak\{\}MS & 5.6625 & 3.30358\\
games, 1000, , complex, 1000 & avl & SEARCH\_\allowbreak\{\}COMPARISONS & 8348 & 0\\
games, 1000, , complex, 1000 & avl & SEARCH\_\allowbreak\{\}MS & 0.3713 & 0.20611\\
games, 1000, , complex, 1000 & avl & TOTAL\_\allowbreak\{\}MS & 31.1343 & 15.4575\\
games, 1000, , complex, 1000 & avl & WRITE\_\allowbreak\{\}MS & 2.334 & 1.1917\\
games, 1000, , complex, 1000 & rubro\_\allowbreak\{\}negra & FETCH\_\allowbreak\{\}MS & 22.7581 & 12.0319\\
games, 1000, , complex, 1000 & rubro\_\allowbreak\{\}negra & LOAD\_\allowbreak\{\}INDEX\_\allowbreak\{\}MS & 5.5574 & 2.6194\\
games, 1000, , complex, 1000 & rubro\_\allowbreak\{\}negra & SEARCH\_\allowbreak\{\}COMPARISONS & 8145 & 0\\
games, 1000, , complex, 1000 & rubro\_\allowbreak\{\}negra & SEARCH\_\allowbreak\{\}MS & 0.3489 & 0.170177\\
games, 1000, , complex, 1000 & rubro\_\allowbreak\{\}negra & TOTAL\_\allowbreak\{\}MS & 32.6038 & 16.5587\\
games, 1000, , complex, 1000 & rubro\_\allowbreak\{\}negra & WRITE\_\allowbreak\{\}MS & 2.5024 & 1.40801\\
games, 1000, , simple, 1 & abb & FETCH\_\allowbreak\{\}MS & 0.0321 & 0.0193187\\
games, 1000, , simple, 1 & abb & LOAD\_\allowbreak\{\}INDEX\_\allowbreak\{\}MS & 1.0606 & 0.555235\\
games, 1000, , simple, 1 & abb & SEARCH\_\allowbreak\{\}COMPARISONS & 9 & 0\\
games, 1000, , simple, 1 & abb & SEARCH\_\allowbreak\{\}MS & 0.0011 & 0.000737865\\
games, 1000, , simple, 1 & abb & TOTAL\_\allowbreak\{\}MS & 1.2283 & 0.652873\\
games, 1000, , simple, 1 & abb & WRITE\_\allowbreak\{\}MS & 0.0226 & 0.0140016\\
games, 1000, , simple, 1 & avl & FETCH\_\allowbreak\{\}MS & 0.0306 & 0.0191845\\
games, 1000, , simple, 1 & avl & LOAD\_\allowbreak\{\}INDEX\_\allowbreak\{\}MS & 1.1276 & 0.702628\\
games, 1000, , simple, 1 & avl & SEARCH\_\allowbreak\{\}COMPARISONS & 8 & 0\\
games, 1000, , simple, 1 & avl & SEARCH\_\allowbreak\{\}MS & 0.0009 & 0.000567646\\
games, 1000, , simple, 1 & avl & TOTAL\_\allowbreak\{\}MS & 1.2892 & 0.795195\\
games, 1000, , simple, 1 & avl & WRITE\_\allowbreak\{\}MS & 0.0215 & 0.0114431\\
games, 1000, , simple, 1 & rubro\_\allowbreak\{\}negra & FETCH\_\allowbreak\{\}MS & 0.0351 & 0.0230046\\
games, 1000, , simple, 1 & rubro\_\allowbreak\{\}negra & LOAD\_\allowbreak\{\}INDEX\_\allowbreak\{\}MS & 1.2186 & 0.802271\\
games, 1000, , simple, 1 & rubro\_\allowbreak\{\}negra & SEARCH\_\allowbreak\{\}COMPARISONS & 8 & 0\\
games, 1000, , simple, 1 & rubro\_\allowbreak\{\}negra & SEARCH\_\allowbreak\{\}MS & 0.0008 & 0.000632456\\
games, 1000, , simple, 1 & rubro\_\allowbreak\{\}negra & TOTAL\_\allowbreak\{\}MS & 1.402 & 0.912689\\
games, 1000, , simple, 1 & rubro\_\allowbreak\{\}negra & WRITE\_\allowbreak\{\}MS & 0.0241 & 0.0137069\\
games, 1000, , simple, 10 & abb & FETCH\_\allowbreak\{\}MS & 0.3763 & 0.224402\\
games, 1000, , simple, 10 & abb & LOAD\_\allowbreak\{\}INDEX\_\allowbreak\{\}MS & 1.26 & 0.751257\\
games, 1000, , simple, 10 & abb & SEARCH\_\allowbreak\{\}COMPARISONS & 100 & 0\\
games, 1000, , simple, 10 & abb & SEARCH\_\allowbreak\{\}MS & 0.0049 & 0.0025144\\
games, 1000, , simple, 10 & abb & TOTAL\_\allowbreak\{\}MS & 1.8424 & 1.09543\\
games, 1000, , simple, 10 & abb & WRITE\_\allowbreak\{\}MS & 0.0619 & 0.0371796\\
games, 1000, , simple, 10 & avl & FETCH\_\allowbreak\{\}MS & 0.4189 & 0.294963\\
games, 1000, , simple, 10 & avl & LOAD\_\allowbreak\{\}INDEX\_\allowbreak\{\}MS & 1.3927 & 0.876998\\
games, 1000, , simple, 10 & avl & SEARCH\_\allowbreak\{\}COMPARISONS & 77 & 0\\
games, 1000, , simple, 10 & avl & SEARCH\_\allowbreak\{\}MS & 0.0052 & 0.00319026\\
games, 1000, , simple, 10 & avl & TOTAL\_\allowbreak\{\}MS & 2.0735 & 1.35158\\
games, 1000, , simple, 10 & avl & WRITE\_\allowbreak\{\}MS & 0.0753 & 0.0561645\\
games, 1000, , simple, 10 & rubro\_\allowbreak\{\}negra & FETCH\_\allowbreak\{\}MS & 0.2964 & 0.129128\\
games, 1000, , simple, 10 & rubro\_\allowbreak\{\}negra & LOAD\_\allowbreak\{\}INDEX\_\allowbreak\{\}MS & 1.1039 & 0.558241\\
games, 1000, , simple, 10 & rubro\_\allowbreak\{\}negra & SEARCH\_\allowbreak\{\}COMPARISONS & 80 & 0\\
games, 1000, , simple, 10 & rubro\_\allowbreak\{\}negra & SEARCH\_\allowbreak\{\}MS & 0.0049 & 0.0025144\\
games, 1000, , simple, 10 & rubro\_\allowbreak\{\}negra & TOTAL\_\allowbreak\{\}MS & 1.5875 & 0.770582\\
games, 1000, , simple, 10 & rubro\_\allowbreak\{\}negra & WRITE\_\allowbreak\{\}MS & 0.0526 & 0.0254611\\
games, 1000, , simple, 100 & abb & FETCH\_\allowbreak\{\}MS & 3.6428 & 1.80424\\
games, 1000, , simple, 100 & abb & LOAD\_\allowbreak\{\}INDEX\_\allowbreak\{\}MS & 1.0823 & 0.461127\\
games, 1000, , simple, 100 & abb & SEARCH\_\allowbreak\{\}COMPARISONS & 938 & 0\\
games, 1000, , simple, 100 & abb & SEARCH\_\allowbreak\{\}MS & 0.0331 & 0.0146776\\
games, 1000, , simple, 100 & abb & TOTAL\_\allowbreak\{\}MS & 5.413 & 2.62146\\
games, 1000, , simple, 100 & abb & WRITE\_\allowbreak\{\}MS & 0.4278 & 0.260842\\
games, 1000, , simple, 100 & avl & FETCH\_\allowbreak\{\}MS & 3.8752 & 2.85458\\
games, 1000, , simple, 100 & avl & LOAD\_\allowbreak\{\}INDEX\_\allowbreak\{\}MS & 1.2123 & 0.812816\\
games, 1000, , simple, 100 & avl & SEARCH\_\allowbreak\{\}COMPARISONS & 759 & 0\\
games, 1000, , simple, 100 & avl & SEARCH\_\allowbreak\{\}MS & 0.03 & 0.0157833\\
games, 1000, , simple, 100 & avl & TOTAL\_\allowbreak\{\}MS & 5.784 & 4.08932\\
games, 1000, , simple, 100 & avl & WRITE\_\allowbreak\{\}MS & 0.4093 & 0.274724\\
games, 1000, , simple, 100 & rubro\_\allowbreak\{\}negra & FETCH\_\allowbreak\{\}MS & 3.9374 & 1.83086\\
games, 1000, , simple, 100 & rubro\_\allowbreak\{\}negra & LOAD\_\allowbreak\{\}INDEX\_\allowbreak\{\}MS & 1.2569 & 0.523691\\
games, 1000, , simple, 100 & rubro\_\allowbreak\{\}negra & SEARCH\_\allowbreak\{\}COMPARISONS & 756 & 0\\
games, 1000, , simple, 100 & rubro\_\allowbreak\{\}negra & SEARCH\_\allowbreak\{\}MS & 0.0301 & 0.0120319\\
games, 1000, , simple, 100 & rubro\_\allowbreak\{\}negra & TOTAL\_\allowbreak\{\}MS & 5.905 & 2.65339\\
games, 1000, , simple, 100 & rubro\_\allowbreak\{\}negra & WRITE\_\allowbreak\{\}MS & 0.4312 & 0.204099\\
games, 1000, , simple, 1000 & abb & FETCH\_\allowbreak\{\}MS & 24.0892 & 12.7059\\
games, 1000, , simple, 1000 & abb & LOAD\_\allowbreak\{\}INDEX\_\allowbreak\{\}MS & 1.2894 & 0.760977\\
games, 1000, , simple, 1000 & abb & SEARCH\_\allowbreak\{\}COMPARISONS & 9407 & 0\\
games, 1000, , simple, 1000 & abb & SEARCH\_\allowbreak\{\}MS & 0.3343 & 0.155987\\
games, 1000, , simple, 1000 & abb & TOTAL\_\allowbreak\{\}MS & 29.6235 & 15.8097\\
games, 1000, , simple, 1000 & abb & WRITE\_\allowbreak\{\}MS & 2.7021 & 1.56733\\
games, 1000, , simple, 1000 & avl & FETCH\_\allowbreak\{\}MS & 24.6722 & 12.5775\\
games, 1000, , simple, 1000 & avl & LOAD\_\allowbreak\{\}INDEX\_\allowbreak\{\}MS & 1.1888 & 0.584687\\
games, 1000, , simple, 1000 & avl & SEARCH\_\allowbreak\{\}COMPARISONS & 7422 & 0\\
games, 1000, , simple, 1000 & avl & SEARCH\_\allowbreak\{\}MS & 0.25 & 0.1061\\
games, 1000, , simple, 1000 & avl & TOTAL\_\allowbreak\{\}MS & 29.9715 & 15.0812\\
games, 1000, , simple, 1000 & avl & WRITE\_\allowbreak\{\}MS & 2.7468 & 1.4668\\
games, 1000, , simple, 1000 & rubro\_\allowbreak\{\}negra & FETCH\_\allowbreak\{\}MS & 21.2515 & 9.81378\\
games, 1000, , simple, 1000 & rubro\_\allowbreak\{\}negra & LOAD\_\allowbreak\{\}INDEX\_\allowbreak\{\}MS & 1.231 & 0.723775\\
games, 1000, , simple, 1000 & rubro\_\allowbreak\{\}negra & SEARCH\_\allowbreak\{\}COMPARISONS & 7461 & 0\\
games, 1000, , simple, 1000 & rubro\_\allowbreak\{\}negra & SEARCH\_\allowbreak\{\}MS & 0.2719 & 0.135035\\
games, 1000, , simple, 1000 & rubro\_\allowbreak\{\}negra & TOTAL\_\allowbreak\{\}MS & 26.2613 & 12.3898\\
games, 1000, , simple, 1000 & rubro\_\allowbreak\{\}negra & WRITE\_\allowbreak\{\}MS & 2.2884 & 1.10717\\
games, 10000, , complex, 10 & abb & FETCH\_\allowbreak\{\}MS & 0.4153 & 0.0571743\\
games, 10000, , complex, 10 & abb & LOAD\_\allowbreak\{\}INDEX\_\allowbreak\{\}MS & 12.0362 & 2.8049\\
games, 10000, , complex, 10 & abb & SEARCH\_\allowbreak\{\}COMPARISONS & 147 & 0\\
games, 10000, , complex, 10 & abb & SEARCH\_\allowbreak\{\}MS & 0.0081 & 0.00128668\\
games, 10000, , complex, 10 & abb & TOTAL\_\allowbreak\{\}MS & 13.1018 & 2.92055\\
games, 10000, , complex, 10 & abb & WRITE\_\allowbreak\{\}MS & 0.0636 & 0.0123306\\
games, 10000, , complex, 10 & avl & FETCH\_\allowbreak\{\}MS & 0.499 & 0.167154\\
games, 10000, , complex, 10 & avl & LOAD\_\allowbreak\{\}INDEX\_\allowbreak\{\}MS & 11.3956 & 2.1212\\
games, 10000, , complex, 10 & avl & SEARCH\_\allowbreak\{\}COMPARISONS & 109 & 0\\
games, 10000, , complex, 10 & avl & SEARCH\_\allowbreak\{\}MS & 0.0063 & 0.00149443\\
games, 10000, , complex, 10 & avl & TOTAL\_\allowbreak\{\}MS & 12.6206 & 2.32237\\
games, 10000, , complex, 10 & avl & WRITE\_\allowbreak\{\}MS & 0.0739 & 0.0298606\\
games, 10000, , complex, 10 & rubro\_\allowbreak\{\}negra & FETCH\_\allowbreak\{\}MS & 0.4861 & 0.15428\\
games, 10000, , complex, 10 & rubro\_\allowbreak\{\}negra & LOAD\_\allowbreak\{\}INDEX\_\allowbreak\{\}MS & 13.102 & 2.7543\\
games, 10000, , complex, 10 & rubro\_\allowbreak\{\}negra & SEARCH\_\allowbreak\{\}COMPARISONS & 107 & 0\\
games, 10000, , complex, 10 & rubro\_\allowbreak\{\}negra & SEARCH\_\allowbreak\{\}MS & 0.0063 & 0.00156702\\
games, 10000, , complex, 10 & rubro\_\allowbreak\{\}negra & TOTAL\_\allowbreak\{\}MS & 14.2818 & 2.95353\\
games, 10000, , complex, 10 & rubro\_\allowbreak\{\}negra & WRITE\_\allowbreak\{\}MS & 0.0649 & 0.0179224\\
games, 10000, , complex, 100 & abb & FETCH\_\allowbreak\{\}MS & 7.2685 & 1.74747\\
games, 10000, , complex, 100 & abb & LOAD\_\allowbreak\{\}INDEX\_\allowbreak\{\}MS & 11.7624 & 2.25109\\
games, 10000, , complex, 100 & abb & SEARCH\_\allowbreak\{\}COMPARISONS & 1386 & 0\\
games, 10000, , complex, 100 & abb & SEARCH\_\allowbreak\{\}MS & 0.0488 & 0.00828385\\
games, 10000, , complex, 100 & abb & TOTAL\_\allowbreak\{\}MS & 20.5975 & 4.02255\\
games, 10000, , complex, 100 & abb & WRITE\_\allowbreak\{\}MS & 0.7587 & 0.196885\\
games, 10000, , complex, 100 & avl & FETCH\_\allowbreak\{\}MS & 7.39 & 1.75286\\
games, 10000, , complex, 100 & avl & LOAD\_\allowbreak\{\}INDEX\_\allowbreak\{\}MS & 11.9844 & 2.61845\\
games, 10000, , complex, 100 & avl & SEARCH\_\allowbreak\{\}COMPARISONS & 1042 & 0\\
games, 10000, , complex, 100 & avl & SEARCH\_\allowbreak\{\}MS & 0.0464 & 0.0144391\\
games, 10000, , complex, 100 & avl & TOTAL\_\allowbreak\{\}MS & 20.9495 & 4.28823\\
games, 10000, , complex, 100 & avl & WRITE\_\allowbreak\{\}MS & 0.7753 & 0.229904\\
games, 10000, , complex, 100 & rubro\_\allowbreak\{\}negra & FETCH\_\allowbreak\{\}MS & 6.797 & 0.832384\\
games, 10000, , complex, 100 & rubro\_\allowbreak\{\}negra & LOAD\_\allowbreak\{\}INDEX\_\allowbreak\{\}MS & 12.6379 & 1.9164\\
games, 10000, , complex, 100 & rubro\_\allowbreak\{\}negra & SEARCH\_\allowbreak\{\}COMPARISONS & 1083 & 0\\
games, 10000, , complex, 100 & rubro\_\allowbreak\{\}negra & SEARCH\_\allowbreak\{\}MS & 0.0404 & 0.0160153\\
games, 10000, , complex, 100 & rubro\_\allowbreak\{\}negra & TOTAL\_\allowbreak\{\}MS & 20.8954 & 2.53165\\
games, 10000, , complex, 100 & rubro\_\allowbreak\{\}negra & WRITE\_\allowbreak\{\}MS & 0.7136 & 0.0930844\\
games, 10000, , complex, 1000 & abb & FETCH\_\allowbreak\{\}MS & 42.6596 & 7.96853\\
games, 10000, , complex, 1000 & abb & LOAD\_\allowbreak\{\}INDEX\_\allowbreak\{\}MS & 11.2287 & 2.07808\\
games, 10000, , complex, 1000 & abb & SEARCH\_\allowbreak\{\}COMPARISONS & 13758 & 0\\
games, 10000, , complex, 1000 & abb & SEARCH\_\allowbreak\{\}MS & 0.374 & 0.077676\\
games, 10000, , complex, 1000 & abb & TOTAL\_\allowbreak\{\}MS & 60.1962 & 10.9649\\
games, 10000, , complex, 1000 & abb & WRITE\_\allowbreak\{\}MS & 4.4001 & 0.847886\\
games, 10000, , complex, 1000 & avl & FETCH\_\allowbreak\{\}MS & 40.2258 & 7.40411\\
games, 10000, , complex, 1000 & avl & LOAD\_\allowbreak\{\}INDEX\_\allowbreak\{\}MS & 12.6167 & 3.08434\\
games, 10000, , complex, 1000 & avl & SEARCH\_\allowbreak\{\}COMPARISONS & 10230 & 0\\
games, 10000, , complex, 1000 & avl & SEARCH\_\allowbreak\{\}MS & 0.3287 & 0.0556757\\
games, 10000, , complex, 1000 & avl & TOTAL\_\allowbreak\{\}MS & 58.8542 & 10.9058\\
games, 10000, , complex, 1000 & avl & WRITE\_\allowbreak\{\}MS & 4.2065 & 0.826288\\
games, 10000, , complex, 1000 & rubro\_\allowbreak\{\}negra & FETCH\_\allowbreak\{\}MS & 40.2619 & 7.38638\\
games, 10000, , complex, 1000 & rubro\_\allowbreak\{\}negra & LOAD\_\allowbreak\{\}INDEX\_\allowbreak\{\}MS & 12.093 & 3.39056\\
games, 10000, , complex, 1000 & rubro\_\allowbreak\{\}negra & SEARCH\_\allowbreak\{\}COMPARISONS & 10613 & 0\\
games, 10000, , complex, 1000 & rubro\_\allowbreak\{\}negra & SEARCH\_\allowbreak\{\}MS & 0.3516 & 0.0757543\\
games, 10000, , complex, 1000 & rubro\_\allowbreak\{\}negra & TOTAL\_\allowbreak\{\}MS & 58.5932 & 11.4728\\
games, 10000, , complex, 1000 & rubro\_\allowbreak\{\}negra & WRITE\_\allowbreak\{\}MS & 4.2491 & 0.794136\\
games, 10000, , simple, 1 & abb & FETCH\_\allowbreak\{\}MS & 0.0197 & 0.00290784\\
games, 10000, , simple, 1 & abb & LOAD\_\allowbreak\{\}INDEX\_\allowbreak\{\}MS & 2 & 0.238913\\
games, 10000, , simple, 1 & abb & SEARCH\_\allowbreak\{\}COMPARISONS & 9 & 0\\
games, 10000, , simple, 1 & abb & SEARCH\_\allowbreak\{\}MS & 0.0001 & 0.000316228\\
games, 10000, , simple, 1 & abb & TOTAL\_\allowbreak\{\}MS & 2.1413 & 0.250029\\
games, 10000, , simple, 1 & abb & WRITE\_\allowbreak\{\}MS & 0.0116 & 0.00227058\\
games, 10000, , simple, 1 & avl & FETCH\_\allowbreak\{\}MS & 0.019 & 0.000942809\\
games, 10000, , simple, 1 & avl & LOAD\_\allowbreak\{\}INDEX\_\allowbreak\{\}MS & 2.0036 & 0.121978\\
games, 10000, , simple, 1 & avl & SEARCH\_\allowbreak\{\}COMPARISONS & 9 & 0\\
games, 10000, , simple, 1 & avl & SEARCH\_\allowbreak\{\}MS & 0.0005 & 0.000527046\\
games, 10000, , simple, 1 & avl & TOTAL\_\allowbreak\{\}MS & 2.1596 & 0.132372\\
games, 10000, , simple, 1 & avl & WRITE\_\allowbreak\{\}MS & 0.0126 & 0.00241293\\
games, 10000, , simple, 1 & rubro\_\allowbreak\{\}negra & FETCH\_\allowbreak\{\}MS & 0.0186 & 0.00107497\\
games, 10000, , simple, 1 & rubro\_\allowbreak\{\}negra & LOAD\_\allowbreak\{\}INDEX\_\allowbreak\{\}MS & 1.9772 & 0.0735433\\
games, 10000, , simple, 1 & rubro\_\allowbreak\{\}negra & SEARCH\_\allowbreak\{\}COMPARISONS & 10 & 0\\
games, 10000, , simple, 1 & rubro\_\allowbreak\{\}negra & SEARCH\_\allowbreak\{\}MS & 0.0002 & 0.000421637\\
games, 10000, , simple, 1 & rubro\_\allowbreak\{\}negra & TOTAL\_\allowbreak\{\}MS & 2.1203 & 0.0838505\\
games, 10000, , simple, 1 & rubro\_\allowbreak\{\}negra & WRITE\_\allowbreak\{\}MS & 0.011 & 0.00176383\\
games, 10000, , simple, 10 & abb & FETCH\_\allowbreak\{\}MS & 4.9803 & 0.205915\\
games, 10000, , simple, 10 & abb & LOAD\_\allowbreak\{\}INDEX\_\allowbreak\{\}MS & 2.0053 & 0.117237\\
games, 10000, , simple, 10 & abb & SEARCH\_\allowbreak\{\}COMPARISONS & 119 & 0\\
games, 10000, , simple, 10 & abb & SEARCH\_\allowbreak\{\}MS & 0.0031 & 0.000316228\\
games, 10000, , simple, 10 & abb & TOTAL\_\allowbreak\{\}MS & 7.6584 & 0.343063\\
games, 10000, , simple, 10 & abb & WRITE\_\allowbreak\{\}MS & 0.5109 & 0.0573051\\
games, 10000, , simple, 10 & avl & FETCH\_\allowbreak\{\}MS & 5.4154 & 0.805517\\
games, 10000, , simple, 10 & avl & LOAD\_\allowbreak\{\}INDEX\_\allowbreak\{\}MS & 1.9809 & 0.0755461\\
games, 10000, , simple, 10 & avl & SEARCH\_\allowbreak\{\}COMPARISONS & 98 & 0\\
games, 10000, , simple, 10 & avl & SEARCH\_\allowbreak\{\}MS & 0.003 & 0\\
games, 10000, , simple, 10 & avl & TOTAL\_\allowbreak\{\}MS & 8.1565 & 0.933433\\
games, 10000, , simple, 10 & avl & WRITE\_\allowbreak\{\}MS & 0.5742 & 0.0909943\\
games, 10000, , simple, 10 & rubro\_\allowbreak\{\}negra & FETCH\_\allowbreak\{\}MS & 5.0541 & 0.261846\\
games, 10000, , simple, 10 & rubro\_\allowbreak\{\}negra & LOAD\_\allowbreak\{\}INDEX\_\allowbreak\{\}MS & 2.1006 & 0.287173\\
games, 10000, , simple, 10 & rubro\_\allowbreak\{\}negra & SEARCH\_\allowbreak\{\}COMPARISONS & 100 & 0\\
games, 10000, , simple, 10 & rubro\_\allowbreak\{\}negra & SEARCH\_\allowbreak\{\}MS & 0.0027 & 0.000674949\\
games, 10000, , simple, 10 & rubro\_\allowbreak\{\}negra & TOTAL\_\allowbreak\{\}MS & 7.8548 & 0.555316\\
games, 10000, , simple, 10 & rubro\_\allowbreak\{\}negra & WRITE\_\allowbreak\{\}MS & 0.5195 & 0.0536185\\
games, 10000, , simple, 100 & abb & FETCH\_\allowbreak\{\}MS & 10.851 & 0.547942\\
games, 10000, , simple, 100 & abb & LOAD\_\allowbreak\{\}INDEX\_\allowbreak\{\}MS & 2.2054 & 0.32133\\
games, 10000, , simple, 100 & abb & SEARCH\_\allowbreak\{\}COMPARISONS & 1091 & 0\\
games, 10000, , simple, 100 & abb & SEARCH\_\allowbreak\{\}MS & 0.0235 & 0.00135401\\
games, 10000, , simple, 100 & abb & TOTAL\_\allowbreak\{\}MS & 14.4603 & 0.776561\\
games, 10000, , simple, 100 & abb & WRITE\_\allowbreak\{\}MS & 1.1088 & 0.0833477\\
games, 10000, , simple, 100 & avl & FETCH\_\allowbreak\{\}MS & 10.8666 & 0.762303\\
games, 10000, , simple, 100 & avl & LOAD\_\allowbreak\{\}INDEX\_\allowbreak\{\}MS & 2.1966 & 0.373057\\
games, 10000, , simple, 100 & avl & SEARCH\_\allowbreak\{\}COMPARISONS & 944 & 0\\
games, 10000, , simple, 100 & avl & SEARCH\_\allowbreak\{\}MS & 0.0268 & 0.0145128\\
games, 10000, , simple, 100 & avl & TOTAL\_\allowbreak\{\}MS & 14.4901 & 1.20004\\
games, 10000, , simple, 100 & avl & WRITE\_\allowbreak\{\}MS & 1.1141 & 0.13266\\
games, 10000, , simple, 100 & rubro\_\allowbreak\{\}negra & FETCH\_\allowbreak\{\}MS & 11.8055 & 1.6165\\
games, 10000, , simple, 100 & rubro\_\allowbreak\{\}negra & LOAD\_\allowbreak\{\}INDEX\_\allowbreak\{\}MS & 2.3716 & 0.434004\\
games, 10000, , simple, 100 & rubro\_\allowbreak\{\}negra & SEARCH\_\allowbreak\{\}COMPARISONS & 934 & 0\\
games, 10000, , simple, 100 & rubro\_\allowbreak\{\}negra & SEARCH\_\allowbreak\{\}MS & 0.0229 & 0.00460555\\
games, 10000, , simple, 100 & rubro\_\allowbreak\{\}negra & TOTAL\_\allowbreak\{\}MS & 15.7636 & 2.07826\\
games, 10000, , simple, 100 & rubro\_\allowbreak\{\}negra & WRITE\_\allowbreak\{\}MS & 1.2421 & 0.226395\\
games, 10000, , simple, 1000 & abb & FETCH\_\allowbreak\{\}MS & 63.3138 & 9.21527\\
games, 10000, , simple, 1000 & abb & LOAD\_\allowbreak\{\}INDEX\_\allowbreak\{\}MS & 2.1759 & 0.161071\\
games, 10000, , simple, 1000 & abb & SEARCH\_\allowbreak\{\}COMPARISONS & 11752 & 0\\
games, 10000, , simple, 1000 & abb & SEARCH\_\allowbreak\{\}MS & 0.2338 & 0.015683\\
games, 10000, , simple, 1000 & abb & TOTAL\_\allowbreak\{\}MS & 73.5445 & 10.335\\
games, 10000, , simple, 1000 & abb & WRITE\_\allowbreak\{\}MS & 6.6228 & 0.935939\\
games, 10000, , simple, 1000 & avl & FETCH\_\allowbreak\{\}MS & 66.2865 & 9.94994\\
games, 10000, , simple, 1000 & avl & LOAD\_\allowbreak\{\}INDEX\_\allowbreak\{\}MS & 2.4334 & 0.391153\\
games, 10000, , simple, 1000 & avl & SEARCH\_\allowbreak\{\}COMPARISONS & 9546 & 0\\
games, 10000, , simple, 1000 & avl & SEARCH\_\allowbreak\{\}MS & 0.2202 & 0.0435655\\
games, 10000, , simple, 1000 & avl & TOTAL\_\allowbreak\{\}MS & 77.2943 & 11.5313\\
games, 10000, , simple, 1000 & avl & WRITE\_\allowbreak\{\}MS & 7.0396 & 1.11508\\
games, 10000, , simple, 1000 & rubro\_\allowbreak\{\}negra & FETCH\_\allowbreak\{\}MS & 62.042 & 7.6138\\
games, 10000, , simple, 1000 & rubro\_\allowbreak\{\}negra & LOAD\_\allowbreak\{\}INDEX\_\allowbreak\{\}MS & 2.7669 & 0.79413\\
games, 10000, , simple, 1000 & rubro\_\allowbreak\{\}negra & SEARCH\_\allowbreak\{\}COMPARISONS & 9327 & 0\\
games, 10000, , simple, 1000 & rubro\_\allowbreak\{\}negra & SEARCH\_\allowbreak\{\}MS & 0.216 & 0.0431303\\
games, 10000, , simple, 1000 & rubro\_\allowbreak\{\}negra & TOTAL\_\allowbreak\{\}MS & 73.2913 & 8.93819\\
games, 10000, , simple, 1000 & rubro\_\allowbreak\{\}negra & WRITE\_\allowbreak\{\}MS & 6.9008 & 1.03844\\
\hline
\end{longtable}
\normalsize
}{\TODO{executar experimentos/analise.py para inserir as tabelas de busca}}

\subsection{Estatística pareada}
Para as diferenças pareadas $d_i=x_i^{(A)}-x_i^{(B)}$, calcula-se $\bar d$, o desvio padrão amostral $s_d$ e $t=\bar d/(s_d/\sqrt{r})$. Para dez repetições, o valor crítico bilateral a 5\% com nove graus de liberdade é 2,262. Rejeita-se $H_0$ quando $|t|>2,262$. A análise automática também registra o valor-p se SciPy estiver instalado.

Na configuração de asteroides, 1000 registros e consulta complexa de 10 chaves, ABB contra AVL produziu $t=2{,}753$, acima do crítico, enquanto AVL contra rubro-negra produziu $t=-1{,}000$. Na configuração de games, 10000 registros e busca simples de 1000 chaves, os valores foram $t=1{,}005$ para ABB contra AVL e $t=0{,}234$ para AVL contra rubro-negra. Assim, somente o primeiro par apresentou diferença estatisticamente detectável entre as quatro comparações selecionadas. A diferença média desse par foi de apenas 0,0008 ms; isso não representa, por si só, um ganho prático relevante. Os tempos de busca são muito pequenos diante dos custos de carregar o índice e buscar/gravar registros, e a impressão com três casas decimais limita a leitura de diferenças muito finas.

\IfFileExists{../experimentos/dados/testes_pareados.tex}{Configuração: asteroids, 1000, complex, 10\\
Par: abb-avl; r=10; t=2.753; t crítico=2.262; rejeita H0=True\\
Diferença média=0.0008 ms; desvio=0.000918937 ms
Configuração: asteroids, 1000, complex, 10\\
Par: avl-rubro\_\allowbreak\{\}negra; r=10; t=-1; t crítico=2.262; rejeita H0=False\\
Diferença média=-0.0002 ms; desvio=0.000632456 ms
Configuração: games, 10000, simple, 1000\\
Par: abb-avl; r=10; t=1.005; t crítico=2.262; rejeita H0=False\\
Diferença média=0.0136 ms; desvio=0.0427738 ms
Configuração: games, 10000, simple, 1000\\
Par: avl-rubro\_\allowbreak\{\}negra; r=10; t=0.2336; t crítico=2.262; rejeita H0=False\\
Diferença média=0.0042 ms; desvio=0.0568601 ms
}{\TODO{arquivo de teste pareado ausente}}

\subsection{Gráficos}
\grafico{construcao_asteroids_year}{Construção por tamanho: asteroides, chave year.}
\grafico{construcao_games_Genres}{Construção por tamanho: games, chave Genres.}
\grafico{busca_simple_asteroids}{Busca simples por quantidade de chaves: asteroides.}
\grafico{busca_complex_asteroids}{Busca complexa por quantidade de chaves: asteroides.}
\grafico{busca_simple_games}{Busca simples por quantidade de chaves: games.}
\grafico{busca_complex_games}{Busca complexa por quantidade de chaves: games.}
\grafico{operacoes_insert_comparisons_asteroids_year}{Comparações de inserção em função do tempo de construção.}
\grafico{operacoes_rotations_asteroids_year}{Rotações em função do tempo de construção.}
\grafico{operacoes_tree_height_games_Genres}{Altura final em função do tempo de construção.}
\grafico{custo_load_index_ms_simple_asteroids}{Tempo de carga do índice nas buscas simples de asteroides.}
\grafico{custo_fetch_ms_simple_asteroids}{Tempo de leitura dos registros nas buscas simples de asteroides.}
\grafico{custo_write_ms_simple_asteroids}{Tempo de escrita nas buscas simples de asteroides.}
\grafico{custo_total_ms_simple_asteroids}{Custo total da busca simples em asteroides.}
\grafico{custo_total_ms_complex_games}{Custo total da busca complexa em games.}

\end{document}
